\section{Positive contact constant and slow-root collision}
\label{sec:positive-constant}

The remaining contact scaling has two independent small parameters:
\begin{equation}
\label{eq:pc-family}
 F=y(1+bz-y)+\varepsilon^2P(z),\qquad
 G=zy+\kappa^3d,\qquad
 P(z)=z^2+arz+r^2,\qquad \kappa=r\varepsilon.
\end{equation}
Throughout this section, $a,d$ range in fixed compact intervals and
$r,\varepsilon>0$ tend to zero independently. In particular, estimates
must allow $r$ to be smaller than every power of $\varepsilon$.

\begin{proposition}
\label{prop:positive-constant}
On every compact $(a,d,b)$-set with $0<b<2$, all sufficiently small
fields \eqref{eq:pc-family} have total actual minimal-return multiplicity
at most two for isolated cycles in one Hausdorff neighborhood of the
whole limiting graphic. The parameter
and Hausdorff neighborhoods are uniform on that compact set.
For $d\ge0$ the bound is one and every counted zero is simple.
\end{proposition}

Near a positive double root, the slow passage can amplify the incoming
height by an arbitrarily large factor. We retain this multiplier in
product coordinates. Weighting the complete outgoing derivative by the
incoming distance from the separating trajectory cancels the multiplier
and gives size and logarithmic-curvature bounds. Combined with the
incoming distortion estimate, these bounds force the derivative ratio
with the short map to cross one only upwards
(Lemma~\ref{lem:pc-ratio}).

Lemma~\ref{lem:pc-entry} controls the incoming trace, and
Lemma~\ref{lem:pc-normal-derivative} estimates the logarithmic derivative
from departure to the target. Before applying the comparison at a
collision, Lemma~\ref{lem:pc-collision-localization} places every counted
cycle in one small source interval. Lemma~\ref{lem:pc-source-partition}
divides it into through passages and fiber exits; on the root side only
the fiber exit occurs (Figure~\ref{fig:collision-passages}). At their
possible finite junction, the two formulas represent the same physical
return. Lemma~\ref{lem:pc-event-gluing} therefore gives the bound two on
the whole interval.

\subsection{The incoming map in a regular coordinate}

We first transport the upper-section source to the inner cut $z=0$.
The signed variation of its height will control the derivative of the
complete contact map.

Use the signed coordinate $u_{\rm up}$ on the upper section from
Lemma~\ref{lem:sections}. Every counted cycle meets this section exactly
once.
The same lemma supplies the contact passage in $y>0$, $F>0$,
from the decreasing-height section $y=K\varepsilon^3$ to a fixed
positive section, on the interval of sources for which that passage is
defined. No sign or scaled-size condition is imposed on $u_{\rm up}$.
Set
\begin{equation}
\label{eq:pc-inner}
 z=\kappa\zeta,\quad y=\kappa^2q,\quad
 q_0(u_{\rm up})=q(0;u_{\rm up}),\qquad
 q_\zeta=\frac{\zeta q+d}{V(\zeta,q)},\quad
 V=\mathsf A+\mathsf Bq-\mathsf Cq^2,
\end{equation}
where
\[
 \mathsf A=1+a\varepsilon\zeta+\varepsilon^2\zeta^2,
 \qquad \mathsf B=1+b\kappa\zeta,\qquad \mathsf C=\kappa^2.
\]
All derivatives below keep the vector field fixed.

\begin{lemma}
\label{lem:pc-entry}
The actual incoming values satisfy $0<q_0\le Q_0$, with $Q_0$
uniform. There is a particular solution $q_c$, independent of
$u_{\rm up}$, such
that $q_0=q_c(0)+w_0$, $w_0>0$. For $d\ge0$,
\begin{equation}
\label{eq:pc-entry-positive}
 q_c(0)\ge0,\qquad
 \frac{|(q_0)_{u_{\rm up}}|}{q_0+d}\le C\varepsilon^{-2}.
\end{equation}
For $d=-m<0$, a cycle is possible only if
\begin{equation}
\label{eq:pc-exponential-wedge}
 a<0,\qquad 0<m\le C\exp(-c/\varepsilon^2).
\end{equation}
On this set $-4m\le q_c(0)<0$ and
\begin{equation}
\label{eq:pc-entry-negative}
 (\log w_0)_{u_{\rm up}}=-\varepsilon^{-2}\lambda,
 \qquad c\le\lambda\le C,\qquad
 |(\log\lambda)_{u_{\rm up}}|\le C\varepsilon^{-2}.
\end{equation}
The passing and target restrictions have interval preimages in
$u_{\rm up}$.
\end{lemma}

\begin{proof}
We first justify the full incoming passage. On a fixed negative interval,
$P(z)\ge z^2/2$ for small $r$. Put $x=-z$ and subtract the forcing
graph $\kappa^3d/x$. For $x\ge Mr$, with $M$ fixed and large, the
two sides of
\[
 \left|y-\frac{\kappa^3d}{x}\right|
       \le k\varepsilon^2x^2
\]
point inward; this follows by substituting the two equalities in
$dy/dz=G/F$, using $P\asymp x^2$ and then choosing $M$ large.
It carries the fixed negative entrance to $Z=z/r=-M$, with bounded
small normal coordinate in
\[
 Y=\frac{y}{r^2\varepsilon^2},\qquad
 w=Y+\frac{\varepsilon d}{Z}.
\]
On a fixed negative $Z$-interval the normal eigenvalue is bounded away
from zero. The invariant graph and strong foliation give an exact
scalar base equation $\dot\theta=\varepsilon^2UQ_\varepsilon(\theta)$,
where $U>0$ and, before the near-identity base change,
\[
 Q_\varepsilon(Z)=Z^2+aZ+1-\frac{\varepsilon d(1+brZ)}Z
                         -\frac{r^2\varepsilon^4d^2}{Z^2}.
\]
Its second derivative is positive on this fixed strip for sufficiently
small parameters, so it has at most two roots. The normal coordinate
contracts at a fixed rate.
The base at entry lies in $(-M-1/2,-M+1/2)$; choosing $M$ larger puts
this entire interval to the left of every possible negative root.
A first root cannot be crossed by the scalar base. Its approach takes
infinite time and normal contraction makes this alternative terminal.
If there is no root, the exit at a fixed $Z=-\zeta_0<0$ satisfies
\[
 |w|\le C\varepsilon^3+C e^{-c/\varepsilon^2}.
\]
The two sides of $|q+d/\zeta|\le K/\zeta^2$ point inward on
$-\zeta_0/\varepsilon\le\zeta\le-L$, and $V$ lies between two
positive constants there. This gives a bounded entrance at $\zeta=-L$.
These assertions use the strong spectral gap, not a lower bound for a
weak root derivative. The fixed-gap graph transform, as in the proof
of Lemma~\ref{lem:small-ratio-saddle}, first supplies the invariant
graph and its finite phase bounds on this negative strip.
Lemma~\ref{lem:uniform-normal-product} then applies with
$\tau=\rho=\varepsilon^2$ and negative normal sign. No scalar base
linearization is made, so its conclusion includes a negative double
root as well.

Here is a direct bound on the last inner ordinate. In $X=z/r$ and
$Y=y/(r^2\varepsilon^2)$ the equations are
\[
 X'=\varepsilon^2(YQ+p_a(X)),\quad
 Y'=XY+\varepsilon d,\quad
 Q=1+brX-r^2\varepsilon^2Y,\quad p_a=X^2+aX+1.
\]
The preceding upper barrier gives $Y(-L_1)\le K_1$ for a fixed large
$L_1$. Choose $T>K_1+4\max_{[-L_1,0]}|p_a|+1$. Above $Y=T$,
$YQ+p_a\ge Y/4$. If $d\le0$, $Y$ decreases for $X<0$. If $d>0$,
an upcrossing of $Y=T$ requires $X\ge-\varepsilon d/T$; thereafter
\[
 \frac{d(Y^2)}{dX}\le\frac{8d}{\varepsilon}.
\]
The remaining interval has length at most $\varepsilon d/T$.
Thus $Y(0)^2\le T^2+8d^2/T$. This proves the uniform bound on $q_0$.

Select the particular solution $y_c$ by zero strong amplitude on the
incoming invariant graph and continue the same solution through the
taper and the inner funnel. The strong-graph equation gives, on every
fixed negative interval,
\[
 y_c=-\kappa^3d/z+O(\varepsilon^2\kappa^3|d|),
 \qquad \partial_z^jy_c=O(\kappa^3|d|),\quad j\le2.
\]
The construction is made on the nested negative intervals with the
same graph, so these are restrictions of one solution. For $d=0$ it
is exactly zero. For $d\ne0$ its sign on the incoming half is the
sign of $d$: at $y=0$ the scalar velocity has that sign, and the
selected forced trace cannot cross to the other side. Its continuation
to $z=0$ has the same sign. At the decreasing-height entrance
$z=\Lambda(u_{\rm up})$,
\[
 W_i:=K\varepsilon^3-y_c(\Lambda)\asymp\varepsilon^3>0.
\]
Scalar uniqueness preserves $W=y-y_c>0$ along every passing interval.

For the needed variational estimates write, exactly,
\begin{equation}
\label{eq:pc-tail-feedback}
 \begin{aligned}
 W_z&=A(z,W)W,&
 A(z,W)&=\int_0^1H_y(z,y_c+tW)\,dt,\\
 L_z&=WA_WL,& L&=(\log W)_{u_{\rm up}},\qquad H=G/F.
 \end{aligned}
\end{equation}
Because the entrance moves, its fixed-phase initial value is
\begin{equation}
\label{eq:pc-moving-entry}
 L_i=\frac{(W_i)_{u_{\rm up}}}{W_i}
       -A(\Lambda,W_i)\Lambda_{u_{\rm up}}.
\end{equation}
The moving entry satisfies the common $C^2$ estimate
\eqref{eq:relaxation-map}, proved by integrating
\eqref{eq:relaxation-error}; its error from the fast-foot map is
$O(\varepsilon^2(1+|\log\varepsilon|))$.
At this section
\[
 A(\Lambda,W_i)
 =\frac{\Lambda}{\varepsilon^2P(\Lambda)}(1+O(\varepsilon)).
\]
It follows that $|L_i|\le C\varepsilon^{-2}$ and, for negative $d$,
$L_i=-\varepsilon^{-2}\lambda_i$, where $\lambda_i$ is a positive
unit with bounded logarithmic $u_{\rm up}$-derivative.

To propagate the initial variation through the shrinking scales, we
bound the feedback integrals. On the fixed negative interval,
\[
 |H_{yy}|\le C\varepsilon^{-4},\quad
 |H_{yyy}|\le C\varepsilon^{-6},\quad
 W\le C\varepsilon^3e^{-c(z-\Lambda)/\varepsilon^2}.
\]
On the taper from $x_g$ to $Mr$,
\[
 \begin{gathered}
 |H_{yy}|\le C\varepsilon^{-4}x^{-3},\qquad
 |H_{yyy}|\le C\varepsilon^{-6}x^{-5},\\
 W(-x)\le W_g(x/x_g)^N,\qquad N=c/\varepsilon^2,\qquad
 W_g\le C\varepsilon^3e^{-c/\varepsilon^2}.
 \end{gathered}
\]
Integrating these powers down to zero bounds the first feedback by
$CW_g/(\varepsilon^2x_g^2)$ and the second by
$CW_g^2/(\varepsilon^4x_g^4)$. On a passing negative root interval,
the exact strong equation $\dot\eta=-b_n(\theta)\eta$,
$b_n\ge c>0$, bounds the first feedback by
$C\int|\eta(0)|e^{-ct}\,dt$. The endpoint changes
$Z=\theta+O(\varepsilon^2\eta)$ have the same bound. To control a
fixed physical section without dividing by a small base speed, write
$\dot\theta=\varepsilon^2v(\theta)$ on a passing branch, where $v>0$.
The entrance speed is a positive unit, $v'$ is bounded, and
\[
 \frac d{dt}\log\frac{|\eta|}{v}
       =-b_n-\varepsilon^2v'\le-c/2.
\]
Consequently $|\eta|/v\le C|\eta(0)|e^{-ct/2}$, which controls
the ratio occurring in the physical-section correction. Indeed, the
product inverse has $Z=\theta+\varepsilon^2\eta L_1$ and
$u=\eta+\eta^2L_2$, with bounded finite jets. Thus $v(Z)$ is
comparable with $v(\theta)$ and the physical speed
$v(Z)+c(Z,u)u$ is a positive unit times $v(Z)$, provided the entrance
amplitude is fixed sufficiently small. For the divided physical
coefficient $a=B/[\varepsilon^2(v+cu)]$, direct differentiation gives
\[
 |u a_u|\,dZ\le C\left(|u|+\frac{|u|}{v(Z)}\right)dt,
 \quad
 |u^2a_{uu}|\,dZ\le C\left(u^2+\frac{u^2}{v(Z)}
                                      +\frac{u^2}{v(Z)^2}\right)dt.
\]
The preceding decay bounds make both integrals finite with a common
bound. The physical feedback equations are differentiated at fixed
$Z$; the product coordinates are only used to estimate their
coefficients. Weak flight time need not be bounded; a base root
reached only at infinite time is terminal and is excluded from this
passing branch.
For negative $d$, the restriction $a<0$ proved below gives
$p_a(Z)\ge1+Z^2$ for $Z<0$. On the fixed negative scaled interval the
scalar equation then has $A\le-c/\varepsilon^2$ and its first two normal
derivatives are $O(\varepsilon^{-2})$, which also bounds the second
feedback.

For $d=-m<0$, a bound proportional to $m$ will localize the later ratio
estimate. At a fixed small negative contact point $z=-x_*$, the selected
incoming graph has
$q_c(-x_*/\kappa)=-\kappa m/x_*+O(\varepsilon^2\kappa m)$.
On the interval from that point to zero, $a<0$ gives $\mathsf A\ge1$.
The graphs $-4m/(1-\zeta)$ and zero are lower and upper barriers.
Indeed, writing $t=-\zeta$, at the lower graph the numerator
$\zeta q-m$ is nonnegative for $t\ge1/3$. For $0\le t\le1/3$,
$V\ge1/2$ bounds the scalar right-hand side below by $-2m$, while the
lower graph has derivative at most $-9m/4$. The initial value is between
the two barriers;
at the upper graph that right-hand side is $-m/\mathsf A<0$. Thus
$-4m\le q_c(0)<0$, with no reciprocal small scale or root gap.

In the funnel set $q=q_c+w$. Its exact divided coefficient is
\[
 a_f=\frac{\mathsf N}{V(\zeta,q_c)V(\zeta,q_c+w)},\quad
 \mathsf N=\zeta\mathsf A-d\mathsf B+
 \mathsf C\{\zeta q_c^2+2dq_c+w(\zeta q_c+d)\}.
\]
On the funnel ending at $\zeta=-L$, with $L$ fixed and large,
$\mathsf N\le-c(|\zeta|+d)$ for $d\ge0$; in the negative
exponential range it is at most $-c|\zeta|$. The speeds are positive
units, and $B_f=(a_f)_w/a_f$ and $(B_f)_w$ are bounded. Since
$w_\zeta=a_fw<0$ and $(a_f)_{ww}=a_f(B_f^2+(B_f)_w)$, the feedback
integrals are bounded by $C\int(-w_\zeta)d\zeta$. The remaining
bounded inner interval $[-L,0]$ has bounded third height derivatives;
the sign of $a_f$ need not stay negative there.
Consequently
\begin{equation}
\label{eq:pc-feedback-mass}
 \int|WA_W|\,dz\le C,\qquad
 \int(|WA_W|+|W^2A_{WW}|)\,dz\le C\quad(d<0).
\end{equation}
Linear changes of phase leave these integrals unchanged. The strong
coordinate calculation above includes the endpoint terms when that
coordinate is used in the negative root interval.

Equations \eqref{eq:pc-tail-feedback}--\eqref{eq:pc-feedback-mass}
preserve $|L|\le C\varepsilon^{-2}$. For $d\ge0$,
$q_0+d\ge w_0$ proves \eqref{eq:pc-entry-positive}. For $d<0$ write
$L(0)=L_i e^{J(u_{\rm up})}$,
$J=\int_{\Lambda(u_{\rm up})}^0WA_W\,dz$. Then
\[
 J_{u_{\rm up}}=-[WA_W]_{\Lambda}\Lambda_{u_{\rm up}}+
              \int L(WA_W+W^2A_{WW})\,dz=O(\varepsilon^{-2}).
\]
The first term is $O(\varepsilon^{-1})$; the second is controlled by
\eqref{eq:pc-feedback-mass}. This proves
\eqref{eq:pc-entry-negative}.

It remains to justify the restriction used on negative $d$.
Lemma~\ref{lem:damping} excludes $a\ge0$. For $d=-m<0$ and $x\ge L_1r$,
the upper barrier and $F\le C\varepsilon^2x^2$ give
\[
 y(-L_1r)\le K\varepsilon^3
       (L_1r/x_g)^{c/\varepsilon^2}.
\]
The height decreases to $z=-\kappa$. Division by $\kappa^2$ loses
only $r^{-2}$: for $N>2$, $r^{-2}(L_1r/x_g)^N$ is increasing in
$r$. Choosing $r_0<x_g/(2L_1)$ gives $q(-1)\le Ce^{-c/\varepsilon^2}$,
uniformly even for arbitrarily small $r$. On $-1\le\zeta\le0$,
$q_\zeta\le-m/2$ as long as $q>0$. Positivity at zero proves
\eqref{eq:pc-exponential-wedge}. Finally, scalar order, $F_y>0$,
and the fixed target interval show that every intervening solution
between two retained ones is retained. Thus the restrictions are
intervals.
\end{proof}

Let $T_-(u_{\rm up})$ be the short map in the energy coordinate $E_b$,
and $T_+(u_{\rm up})$ the map obtained by the contact passage and outgoing arc.
Lemma~\ref{lem:sections} gives
\begin{equation}
\label{eq:pc-short}
 (T_-)_{u_{\rm up}}=-j(u_{\rm up}),\qquad c\le j\le C,\qquad
 |(\log j)_{u_{\rm up}}|\le C.
\end{equation}
Their intersections correspond to captured periodic orbits; their isolated
intersections are precisely the limit cycles being counted.

\paragraph{Actual-domain scope on the nonnegative side.}
For $a>0$ the negative scaled strip may have a base root. The incoming
product in Lemma~\ref{lem:pc-entry} then cannot pass its first root:
its base approaches only in infinite time while the normal deviation
contracts. Such a field has no retained passage to $z=0$, also at
$d=0$. On a no-root branch that lemma controls feedback on the whole
returning interval.
The outgoing and short-map constants can be fixed on fixed
enlarged physical cores before shrinking $\xi$; regular energy
holonomy has a uniform unit bound there. After these constants
are fixed choose $\xi$, then $r$, and finally $\varepsilon$.

\subsection{Nonnegative forcing and a comparison criterion}

For nonnegative forcing, a weighted variational identity makes the
complete contact derivative smaller than the short-map derivative.

For $d\ge0$, the required outgoing estimate holds even when $P$ has
positive roots. Let $J=\partial_{q_0}q$ in \eqref{eq:pc-inner} and
$h=q_0+d>0$. On $0\le\zeta\le1$, direct integration gives
$q(1)\ge ch$ and $0<J(1)\le C$. For $\zeta\ge1$ put
\[
 \mathcal R=hJ\frac{V(\zeta,q)}q.
\]
The exact cancellation
\begin{equation}
\label{eq:pc-positive-cancellation}
 (\log\mathcal R)'=\frac{(V_\zeta)|_q-d/q}{V}
\end{equation}
is obtained by differentiating the scalar and variational equations.
The function $\Phi=\mathsf Bq-\mathsf Cq^2$ increases, since
$q'\ge0$ and $\mathsf B-2\mathsf Cq=F_y>0$. The integral of
$\mathsf B'q/V$ is bounded: outside a negative well of $\mathsf A$,
or when $q$ is large, $q/V\le C$; on the remaining bounded-height
part use
\[
 \frac{\mathsf B'q}{V}\,d\zeta
 =\frac{\mathsf B'q}{\zeta q+d}\,dq
 \le\frac{\mathsf B'}\zeta\,dq.
\]
There $\zeta\ge c/\varepsilon$. To integrate
$\mathsf A'/(\mathsf A+\Phi)$, pair the decreasing and increasing
branches of the convex quadratic at equal values of $\mathsf A$.
The later denominator is larger, so each paired contribution is
nonpositive, including when $\mathsf A\le0$. Only a final increase
above $\mathsf A(1)>c$ remains. At a fixed section $z=\xi$ this gives
\[
 \mathcal R_{\rm out}\le C(1+\xi^2/r^2),\qquad
 (q_0+d)\partial_{q_0}T_+
          \le C(\kappa^2+\varepsilon^2\xi^2).
\]
The factor $\kappa^2$ converts $q$ to the physical ordinate. The
regular continuation in the $E_b$ coordinate has derivative tending
to one after $\xi$ is fixed. With \eqref{eq:pc-entry-positive},
\[
 |(T_+)_{u_{\rm up}}|\le C(r^2+\xi^2).
\]
Choose $\xi$ small and then $r$ small so that the last bound is less
than $j/2$. The displacement is strictly increasing and has at most one
zero, necessarily simple, on its returning interval.

For negative $d$, we use Lemma~\ref{lem:pc-ratio} after identifying
the physical separator and estimating the departure map.

\subsection{Changing the incoming reference preserves the physical separator}
Write $d=-m<0$ and $u=u_{\rm up}$. Lemma~\ref{lem:pc-entry}
has already placed every counted cycle in the wedge $a<0$ with $m$
exponentially small. To estimate incoming distortion, we may replace
the selected incoming reference $y_c^o$ by the solution $y_c^n$ whose
value at a fixed negative point is zero. The outgoing separator is
still the trace $\beta=y_\sigma(0)/\kappa^2$, with source distance
$s=q_0-\beta>0$. We verify that this change of reference preserves
the complete derivative comparison.

Both incoming solutions obey the barriers
$-4m/(1-\zeta)\le q_c^i(\zeta)<0$, $i=o,n$, on their
common terminal negative interval. For $\zeta\in[-1/8,0]$,
$1\le\mathsf A\le2$, $\mathsf B>0$, and
$1/2\le V(\zeta,q_c^i)\le2$ after shrinking.
Moreover $\zeta q_c^i\le m/2$. Their exact scalar equations
therefore give $(q_c^i)_\zeta\le-m/4$, hence
\[
 -4m\le c_i:=q_c^i(0)\le-m/32 .
\]
The old incoming graph's initial value lies between these same
barriers by its fixed-negative-point estimate in
Lemma~\ref{lem:pc-entry}. No equality of the graphs is assumed.

For the same fixed field and physical incoming value put
$w_i=q_0-c_i=s+w_{\beta,i}$,
$w_{\beta,i}=\beta-c_i$, and
$\lambda_i=-\varepsilon^2\partial_u\log w_i$.
The endpoint bounds and $q_0=\beta+s\ge0$ give
\[
 \frac1{128}\le\frac{w_o}{w_n}\le128,\qquad
 \lambda_n=\lambda_o\,\frac{w_o}{w_n},\qquad
 \lambda_iw_i=-\varepsilon^2s_u .
\]
Differentiating the exact identities, with $c_i$ independent of $u$,
gives
\begin{equation}\label{eq:incoming-reference-cancellation}
 (\log\lambda_n)_u
 =(\log\lambda_o)_u+
                    \varepsilon^{-2}(\lambda_n-\lambda_o).
\end{equation}
The feedback calculation
\eqref{eq:pc-tail-feedback}--\eqref{eq:pc-feedback-mass} applies to the
zero-initial reference as well: it has the same $O(m)$ barriers and
fixed-foot, taper and funnel bounds. Its moving-entry value is
$L_i=-H(\Lambda,K\varepsilon^3)\Lambda_u/
(K\varepsilon^3-y_c^n(\Lambda))$, including the entry motion.
The taper integrals have denominators $N_0-2$ and $2N_0-4$,
$N_0\asymp\varepsilon^{-2}$, so no $\log(1/r)$ or reciprocal
radial loss is introduced. Thus the incoming constants are fixed
before the outgoing sections.

If $c_n^\lambda\le\lambda_n\le C_n^\lambda$ and
$|(\log\lambda_n)_u|\le L_n\varepsilon^{-2}$, the identities
give
\[
 c_n^\lambda/128\le\lambda_o\le128C_n^\lambda,\qquad
 |(\log\lambda_o)_u|
       \le(L_n+129C_n^\lambda)\varepsilon^{-2}.
\]
These bounds hold for arbitrarily small $s>0$. The outgoing map
$T_+(s)$, its derivative $M=\partial_sT_+$, the departure derivative $N$,
and the collision amplitude $A(s)$ still use the same physical source
$\kappa^2(\beta+s)$, selected outgoing graph, and target. The derivative
ratio is therefore
\[
 R=\frac{M|s_u|}{j}
  =\frac{M\varepsilon^{-2}\lambda_ow_o}{j}
  =\frac{M\varepsilon^{-2}\lambda_nw_n}{j}.
\]
Equation~\eqref{eq:incoming-reference-cancellation} also gives
$(\log\lambda_n)_u-\varepsilon^{-2}\lambda_n
 =(\log\lambda_o)_u-\varepsilon^{-2}\lambda_o$.
Thus both references give the same logarithmic derivative of $R$,
including its large terms. The outgoing distance is always
$W_{\rm out}(0)=\kappa^2s$, while $\kappa^2w_n$ is used only to estimate
incoming distortion. The separator and finite-event comparison are
therefore preserved.


\subsection{Departure from the slow region}

We next estimate the derivative from a moving departure point to the
fixed outgoing target. The logarithmic derivative of this map, combined
with the departure equation, will give the curvature required by
Lemma~\ref{lem:pc-ratio}.

Write $d=-m$ and $\nu=\kappa^3m$. First suppose that
\begin{equation}
\label{eq:pc-separated-polynomial}
 P(z)\asymp r^2+z^2\quad(z\ge0),\qquad -2+\delta_a\le a<0.
\end{equation}
There is a unique positive solution $y_\sigma$ of $y_z=H=(zy-\nu)/F$
on this side with
\begin{equation}
\label{eq:pc-boundary-tail}
 y_\sigma(z)=\frac{\nu}{(1+\varepsilon^2)z}+O(\nu z^{-2})
 \quad(z\longrightarrow\infty).
\end{equation}
To construct it, use the backward Green kernel
\[
 K(\zeta,t)=\exp\left(-\int_\zeta^t\frac{x}{\mathsf A(x)}\,dx\right).
\]
Writing $\mathcal N(q)=(\zeta q-m)(V^{-1}-\mathsf A^{-1})$,
the integral equation is
$q(\zeta)=m\int_\zeta^\infty K(\zeta,t)/\mathsf A(t)\,dt
          -\int_\zeta^\infty K(\zeta,t)\mathcal N(q)(t)\,dt$.
The comparison $\mathsf A\asymp1+\varepsilon^2\zeta^2$ gives
$\int_\zeta^\infty K/\mathsf A\le C/(1+\zeta)$, by splitting at
$1$ and $\varepsilon^{-1}$. On the ball
$\sup(1+\zeta)|q|\le Cm$, the second integral has Lipschitz norm
$O(m)$ in this weighted space: use
$|V-\mathsf A|\le C(1+\kappa\zeta)|q|+\kappa^2q^2$ in the same
two ranges. Contraction therefore gives the solution. The coefficient
$1/(1+\varepsilon^2)$ is the balance at infinity. A first zero would
cross downward and could not lead to the positive terminal behavior.
Likewise, at a zero of $zy_\sigma-\nu$ its derivative with respect
to $z$ is $y_\sigma>0$. Hence $y_\sigma>0$ and
$zy_\sigma-\nu<0$ throughout. Scalar order shows that the value
$\beta=y_\sigma(0)/\kappa^2$ separates the outgoing solutions from
those hitting $y=0$. The threshold itself does not reach the outgoing
fast part of the fixed section.
For its size, write $\sigma(\zeta)=y_\sigma(\kappa\zeta)/\kappa^2$.
The strict sign above gives $0<\sigma(1)<m$. On $0\le\zeta\le1$,
$\mathsf A\ge3/4$ after shrinking $\varepsilon$. The graphs zero and
$(3-2\zeta)m$ are reverse-time lower and upper barriers: at the upper
graph the scalar right-hand side is at least $-2m$, its derivative is
$-2m$, and the denominator stays above $1/2$ for small $m$.
Therefore $0<\beta\le3m$, and Lemma~\ref{lem:pc-entry} gives
$0<w_\beta=\beta-q_c(0)\le7m$.

Fix two positive sections $\xi_{\rm in}<\xi_{\rm out}$. Choose their
ratio sufficiently large, then their physical positions, and finally the
parameter neighborhood. Retained cycles meet the fast part of the inner
section and continue through the outer section to the common target.
For each pair fixed in this order, we will bound the logarithmic
derivative of the outgoing continuation uniformly in the final
parameter neighborhood; its right-tail version will also apply near
a double root.
For a small fixed $\theta>0$, let $N(z)>0$ be the derivative with
respect to vertical height of this continuation from
\[
 y=y_\sigma(z)+\theta S(z),\qquad S=\varepsilon^2P.
\]

\paragraph{The outgoing logarithmic derivative.}
The key estimate comes from differentiating the exact energy defect as
the starting point moves along the departure curve. The first assertion
below remains under assumption~\eqref{eq:pc-separated-polynomial}; the
last assertion requires that comparison only on its stated right tail.
In both cases the continuation starts at $z_0$ and ends at the common
target, with the field fixed. This right-tail estimate is what the
collision through passage will use.

\begin{lemma}
\label{lem:pc-normal-derivative}
Assume $d=-m<0$ lies in the exponential wedge
\eqref{eq:pc-exponential-wedge} supplied by Lemma~\ref{lem:pc-entry}.
Fix a sufficiently small $\theta>0$.  Choose the ratio
\(R_\xi=\xi_{\rm out}/\xi_{\rm in}>1\) first (large enough for the
post-hit contraction below), then choose the two physical sections, and
only afterward choose the parameter neighborhood.  With these choices
there is an $L>0$ such that, for every admissible parameter in that
neighborhood and every starting point
\[
 z_0\in[L\kappa,\xi_{\rm in}],\qquad
 y_0=y_\sigma(z_0)+\theta S(z_0),
\]
the complete continuation from the inner section through the outer
section to the common target is defined and its vertical derivative
\(N(z_0)>0\) satisfies
\begin{equation}
\label{eq:pc-normal-bound}
 |z_0\,\partial_{z_0}\log N(z_0)|\le C,
 \qquad L\kappa\le z_0\le\xi_{\rm in}.
\end{equation}
Here $\partial_{z_0}$ differentiates the dependence on the
starting-section coordinate $z_0$ while the vector field, normalized
parameters, section positions, and target section are held fixed.  In
the proof, derivatives in the logarithmic phase variable are likewise
field-fixed; whenever $t=\log\lambda$ is used, the radial derivative is
taken at fixed $\lambda$ (fixed $t$), as indicated by the subscript in
$(z\partial_z)_\lambda$.  The same estimate holds with this starting
interval replaced by $[z_*,\xi_{\rm in}]$ whenever $z_*\ge rZ_b$ and
the interval lies on that right tail (in particular, take
$z_*=rZ_b$ when nonempty), where $Z_b$ ranges in a fixed compact
interval lying to the right of all possible positive roots and
\eqref{eq:pc-separated-polynomial} is required only on that tail.  The
constant in this right-tail assertion depends on the fixed tail and
section data, but not on a positive lower bound for $P$ (or $F$) in the
collision strip to the left of the tail.
\end{lemma}

\begin{proof}
\medskip\noindent\textit{Boundary graph.}\quad
We estimate the flight from its moving starting point through the fixed
outer section to the common target; the two distinct physical sections
leave room for departure and subsequent contraction. Put $x=\log z$,
$Z=z/r$, and $y_\sigma=\nu v/z$. For
\[
 \rho=\varepsilon^2\frac{p_a(Z)}{Z^2},\qquad
 \omega=\frac{\nu v(1+bz-\nu v/z)}{z^3},\qquad w=1-v,
\]
the boundary equation becomes
\[
 w_x=(1+c)w-1,\qquad c=(\rho+\omega)^{-1}.
\]
Here $0<v<1$, $\rho\asymp\varepsilon^2(1+Z^{-2})$ and
the first two Euler derivatives of $\rho$ are $O(\rho)$.
Also $\nu/(zS)$ and its first two Euler derivatives are $O(m)$ for
$Z\ge L\varepsilon$. Although $1+bz$ is unbounded at infinity,
its product with $\nu/(zS)$ is still $O(m)$.
The bounded solution has the exact representation
\[
 w(x)=\int_x^\infty
       \exp\left(-\int_x^t(1+c(s,w(s)))\,ds\right)dt.
\]
At fixed $w$, $|\partial_x\log c|\le C$ and
$|\partial_w\log c|\le Cm$. Since initially
$|w_x|\le1+c\le C\varepsilon^{-2}$, the same bound for
$|d\log c/dx|$ holds along the solution by
\eqref{eq:pc-exponential-wedge}. Splitting the integral after a fixed
short interval gives $|w|\le C/(1+c)$. Differentiating the equation,
\[
 (w_x)_x=(1+c+c_ww)w_x+c_xw,
\]
the first coefficient is at least $1+c/2$ and the last term is bounded.
The same backward kernel gives $|w_x|\le C/(1+c)$, and the equation
then gives $|w_{xx}|\le C$. We have proved
\begin{equation}
\label{eq:pc-boundary-euler}
 \sum_{j=0}^2\frac{|(z\partial_z)^jy_\sigma|}{S}
                   +\frac\nu{zS}\le Cm.
\end{equation}

\medskip\noindent\textit{Energy defect before departure.}\quad
Let $H_0=z/(1+bz-y)$. The fast first integral of
Lemma~\ref{lem:fast-energy} gives the exact cancellation, for
$J_z=H_yJ$,
\begin{equation}
\label{eq:pc-energy-defect}
 \begin{gathered}
 \frac d{dz}\log((E_b)_yJ)=\mathcal C,\qquad
 \mathcal C=\left(\partial_y+\frac{(E_b)_{yy}}{(E_b)_y}\right)(H-H_0),\\
 H-H_0=-\frac{\nu h+zS}{Fh},\qquad h=1+bz-y.
 \end{gathered}
\end{equation}
We estimate the change of this integral when its starting point moves.
Set $W=y-y_\sigma$, $R=W/S$, $\lambda=W/P$, and
$D=(\log R)_z$.
Write $y_0(z_0)=y_\sigma(z_0)+\theta S(z_0)$ and let
$y(z;z_0)$ be the solution starting there. With $\ell_0=\log z_0$,
the derivative $\mathcal V=\partial_{\ell_0}y(z;z_0)$ at fixed $z$
satisfies
\[
 \mathcal V_z=H_y\mathcal V,\qquad
 \mathcal V(z_0)=z_0(y_0'-H(z_0,y_0))=-z_0\theta S(z_0)D(z_0).
\]
Choose $c_f<1/4$ and then decrease $\varepsilon_0$ so that
$\varepsilon^2+c_f<1/3$. Before the first hit $\lambda=c_f$, exact
subtraction gives, with $A_*=-a\ge0$,
\[
 \mathfrak a(z,W)=
 \frac{zF_\sigma-G_\sigma(1+bz-2y_\sigma-W)}
      {F_\sigma F_{\sigma+W}},\qquad
 D=\mathfrak a-\frac{P'}P,
\]
where $F_\sigma=F(z,y_\sigma)$,
$F_{\sigma+W}=F(z,y_\sigma+W)$, and
$G_\sigma=zy_\sigma-\nu$.  Equation
\eqref{eq:pc-boundary-euler}, $z\ge Lr\varepsilon$, and
$W=\lambda P$ imply, at fixed $\lambda$,
\[
 F_\sigma=S(1+e_\sigma),\qquad
 F_{\sigma+W}=P(\varepsilon^2+\lambda)
       (1+e_+),
\]
where
\[
 |e_\sigma|+|z\partial_ze_\sigma|\le Cm,\qquad
 |e_+|+|z\partial_ze_+|\le C(z+W+m).
\]
The quotient of
$G_\sigma(1+bz-2y_\sigma-W)$ by $zF_\sigma$, and its fixed-$\lambda$
Euler derivative, are likewise $O(m)$.
The corresponding principal expression is
\[
 D_0=\frac{z}{P(\varepsilon^2+\lambda)}-\frac{P'}P,\qquad
 zD_0=\frac{Z\{A_*v_*+(1-2v_*)Z\}}
                 {v_*p_a(Z)},\qquad v_*=\varepsilon^2+\lambda.
\]
Since $v_*<1/3$ and
$A_*v_*+(1-2v_*)Z\ge A_*v_*+Z/3$, the principal difference cannot
be lost through cancellation.  Dividing the exact error by $D_0$
therefore yields
\begin{equation}
\label{eq:pc-growing-ratio}
 zD=\frac{Z\{A_*v_*+(1-2v_*)Z\}}
                {v_*p_a(Z)}(1+e_1),\qquad
 v_*=\varepsilon^2+\lambda,\quad
 |e_1|+|\partial_xe_1|\le C(z+W+m).
\end{equation}
The Euler derivative holds at fixed $\lambda$. In particular
$D>0$, $D\ge cz/[P(\varepsilon^2+\lambda)]$, and
$|z\partial_z\log D|\le C$. Division of
\eqref{eq:pc-energy-defect}, using $\partial_y=S^{-1}\partial_R$,
gives, for $j=0,1$ at fixed $R$,
\begin{equation}
\label{eq:pc-defect-quotient}
 \left|(z\partial_z)^j\frac{\mathcal C}{D}\right|
             \le C((1+R)^{-1}+S+m).
\end{equation}
For completeness, $H-H_0=-z(1+R)^{-1}\Phi_0$. The function $\Phi_0$
has bounded derivatives in each of its four arguments
\[
 (1+R)^{-1},\qquad h,\qquad y_\sigma/S,\qquad \nu/(zS).
\]
Vertical differentiation therefore bounds $|\mathcal C|$ by
$Cz(1+R)^{-1}\{[S(1+R)]^{-1}+1\}$, while
$D\ge cz/[S(1+R)]$.  Division gives
$C\{(1+R)^{-1}+S\}$; derivatives of the two boundary ratios contribute
the displayed additional $Cm$.  The same calculation after one fixed-$R$
Euler derivative proves the second estimate in
\eqref{eq:pc-defect-quotient}.

\medskip\noindent\textit{The moving departure point.}\quad
To control the moving starting point, use $t=\log\lambda$. The
endpoints are $t_0=\log(\theta\varepsilon^2)$ and $t_f=\log c_f$,
independent of $z_0$. The equation for $x$ is $x_t=B=(zD)^{-1}$;
thus $\mathcal X=\partial_{\ell_0}x$ satisfies
$\mathcal X_t=B_x\mathcal X$, $\mathcal X(t_0)=1$.
Here $B_x$ is taken at fixed $t$. With $e_1=0$,
\[
 B_0=\frac{v_*p_a(Z)}{Z\{A_*v_*+(1-2v_*)Z\}},\quad
 (B_0)_x=
 \frac{v_*\{A_*Z^2(1-v_*)-A_*v_*-2Z(1-2v_*)\}}
 {Z\{A_*v_*+(1-2v_*)Z\}^2}.
\]
If the numerator is positive, $A_*Z\ge c$ and
$((B_0)_x)_+\le Cv_*$. Equation \eqref{eq:pc-growing-ratio} thus
implies
\[
 (B_x)_+\le Cv_*+C(z+W+m)B.
\]
Both terms have bounded integral in $t$. Indeed
$\int v_*dt=\varepsilon^2\log[c_f/(\theta\varepsilon^2)]+O(c_f)$;
also integration of the lower bound for $D$ gives
$z_f\le C(r+z_0)$, $z_f/z_0\le C/\varepsilon$ at the first hit
$\lambda=c_f$. Since $y_z\le Cz$ and $y_\sigma>0$,
$W(z)\le C\varepsilon^2P(z_0)+Cz^2$. In particular
\[
 \int_{z_0}^{z_f}\frac{z+W+m}{z}\,dz
 \le C\{z_f+z_f^2+
       [\varepsilon^2P(z_0)+m]\log(C/\varepsilon)\}\le C.
\]
This uses \eqref{eq:pc-exponential-wedge}, but no comparison of $r$
with $\varepsilon$. Consequently $0<\mathcal X\le C$.
The hit occurs before $C_*(r+z_0)$: otherwise the integral of $z/P$
would exceed the available finite integral of
$(\varepsilon^2+\lambda)d\lambda/\lambda$.  More explicitly,
$D=(\log\lambda)_z$ and the lower bound for $D$ give
\[
 c\int_{z_0}^{z_f}\frac{z\,dz}{P(z)}
 \le\int_{\theta\varepsilon^2}^{c_f}
           \frac{\varepsilon^2+\lambda}{\lambda}\,d\lambda
 \le C.
\]
Since $P(z)\asymp r^2+z^2$ on this side, the left member controls
$\log[(r+z_f)/(r+z_0)]$, proving the asserted bound with a constant
independent of both small parameters.
Choosing $\xi_{\rm out}/\xi_{\rm in}>2C_*$ and then shrinking $r$
puts the hit before the outer section. For the energy defect before
the hit, the physical-coordinate Leibniz formula is exactly
\[
 \partial_{\ell_0}\int_{z_0}^{z_f}\mathcal C(z,y(z;z_0))\,dz
 =\mathcal C_f\partial_{\ell_0}z_f-z_0\mathcal C_0
       +\int_{z_0}^{z_f}\mathcal C_y\mathcal V\,dz.
\]
Changing to $t$ incorporates both endpoint terms and the variational
integral into
\[
 \partial_{\ell_0}\int_{t_0}^{t_f}\frac{\mathcal C}{D}\,dt
 =\int_{t_0}^{t_f}\mathcal X
       (z\partial_z)_{\lambda}\left(\frac{\mathcal C}{D}\right)dt.
\]
Indeed, along a fixed $t$ level,
$\mathcal V=-zWD\mathcal X$. By \eqref{eq:pc-defect-quotient}, the
absolute value of the last integral is at most
\[
 C\int_{t_0}^{t_f}[(1+R)^{-1}+S+m]dt
 \le C\log(1+\theta^{-1})+
 C[\varepsilon^2(r^2+\xi_{\rm out}^2)+m]
       \log\frac{c_f}{\theta\varepsilon^2}\le C.
\]
We have therefore proved
\[
 |\partial_{\log z_0}\log z_f|\le C,\qquad
 \left|\partial_{\log z_0}\int_{z_0}^{z_f}\mathcal C\,dz\right|\le C.
\]

\medskip\noindent\textit{After departure.}\quad
At this hit $z_f\ge c_r r$. Otherwise its height is at least
$cc_fr^2$, whereas $y_z\le Cz$ bounds it by
$C\varepsilon^2r^2+Cz_f^2$, a contradiction. Now set
\[
 \lambda=\frac yP,\qquad \pi=\frac P{z^2}=\frac{p_a(Z)}{Z^2},
 \qquad \delta_Z=\frac\nu{z^3}=m\varepsilon^3Z^{-3},
 \qquad \tau=\log(z/z_f).
\]
Its initial value is
$\lambda(0)=c_f+y_\sigma(z_f)/P(z_f)$, whose Euler derivative is
$O(m\varepsilon^2)$ by \eqref{eq:pc-boundary-euler}.  Direct division of
$y_z=(zy-\nu)/F$ gives the exact equation
\begin{equation}\label{eq:pc-post-exact}
 \lambda_\tau=
 \frac{\lambda-\delta_Z/\pi}{\pi(\lambda h+\varepsilon^2)}
 -\lambda\frac{Zp_a'(Z)}{p_a(Z)},\qquad
 h=1+bz-\lambda P.
\end{equation}
We first control the transition from bounded $Z$ to a stable strip.
Fix \(\lambda_-=c_f/2\) and \(\lambda_+=3/4\), with the earlier choice
\(c_f<1/4\). After shrinking the parameter neighborhood, the hit value
satisfies
\[
 \lambda(0)=c_f+O(m\varepsilon^2)\in(\lambda_-,\lambda_+).
\]
At this hit \(G=zy-\nu>0\), since
\(G=zP\lambda-\nu=zP(\lambda-\delta_Z/\pi)\) and
\(\delta_Z/\pi\ll c_f\).  If \(G\) had a first zero afterward, then
at that point \(G_z=y+zy_z=y>0\), a contradiction; hence
\(y_z=G/F>0\) as long as the orbit stays in the outgoing phase box.
A small fixed outer section closes the bootstrap $h\in[1/2,3/2]$:
$y_z\le2z$ gives $y\le C(r^2+z^2)$, hence $0<\lambda\le B$.
Decrease that section and then $r$ to obtain
$h\in[1-\eta,1+\eta]$, with $\eta<1/10$.
On $Z\ge c_r$, both $\pi$ and $Zp_a'/p_a$ are bounded and $\pi$ is a
positive unit. Choose
$\lambda_*<\min(\lambda_-,c_f/2)$ so small that
\eqref{eq:pc-post-exact} points upward there, and then require
$\varepsilon^2+\delta_Z/\pi\ll\lambda_*$. Thus
$\lambda\in[\lambda_*,B]$. For a fixed sufficiently large $Q$,
$\pi\to1$ and $Zp_a'/p_a\to2$ give
\[
 \lambda_\tau\ge c_1>0\quad(\lambda_*\le\lambda\le\lambda_-),\qquad
 \lambda_\tau\le-c_1<0\quad(\lambda_+\le\lambda\le B),
 \qquad Z\ge Q.
\]
The band $[\lambda_-,\lambda_+]$ is therefore invariant once $Z\ge Q$
and is reached within a bounded additional time. Let $\tau_Q$ be the
first time both conditions hold, or the outer-section time $\tau_o$
if that occurs first. Its preceding interval has length at most
$\max(0,\log(Q/c_r))+B/c_1$ and common coefficient, defect and
variation bounds on $\lambda\in[\lambda_*,B]$.

For the subsequent variation estimate, separate the large-$Z$ equation
into its limiting part and an integrable remainder. Put
\[
 f_\varepsilon(\lambda)=\frac{\lambda}{\lambda+\varepsilon^2}-2\lambda,
\]
and define $\mathcal R_1$ to be the difference between the right-hand
side of \eqref{eq:pc-post-exact} and $f_\varepsilon$.  Equivalently,
\[
 \mathcal R_1=
 \frac{\lambda-\delta_Z/\pi}{\pi(\lambda h+\varepsilon^2)}
 -\frac{\lambda}{\lambda+\varepsilon^2}
 -\lambda\left(\frac{Zp_a'}{p_a}-2\right).
\]
For the energy defect put $e_b=(E_b)_{yy}/(E_b)_y$.  In these variables
the exact expression in \eqref{eq:pc-energy-defect} is
\[
 \mathcal K:=H-H_0
 =-z\frac{\varepsilon^2\pi+\delta_Zh}
             {\pi h(\lambda h+\varepsilon^2)},\qquad
 \mathcal C=(P^{-1}\partial_\lambda+e_b)\mathcal K.
\]
Thus, with $\mathcal E=Z\partial_Z+z\partial_z$ acting at fixed
$\lambda$, define $\mathcal R_2$ by
\begin{equation}
\label{eq:pc-post-remainders}
 \lambda_\tau=f_\varepsilon(\lambda)+\mathcal R_1,\qquad
 z\mathcal C=\frac{\varepsilon^2}{(\lambda+\varepsilon^2)^2}
          +\mathcal R_2.
\end{equation}
On $\lambda_-\le\lambda\le\lambda_+$ and $Z\ge Q$, the denominator
factors are positive units.  The three elementary estimates
\[
 \pi-1=O_{\mathcal E}(Z^{-1}),\qquad
 \frac{Zp_a'}{p_a}-2=O_{\mathcal E}(Z^{-1}),\qquad
 h-1=O_{\mathcal E,\partial_\lambda}(z)
\]
follow directly from the exact formulas
\[
 \pi-1=\frac aZ+\frac1{Z^2},\qquad
 \frac{Zp_a'}{p_a}-2=-\frac{aZ+2}{p_a(Z)},\qquad
 h-1=bz-\lambda\pi z^2,\qquad
 \frac{\delta_Z}{\pi}=\frac{m\varepsilon^3}{Zp_a(Z)}.
\]
The function $e_b$ and its required derivatives are bounded because
$(E_b)_y$ is a positive analytic unit on the fixed energy rectangle.
Substitution in the displayed exact formulas for $\mathcal R_1$ and
$\mathcal K$, followed by the two operations
$P^{-1}\partial_\lambda$ and multiplication by $z$, gives, including
all mixed derivatives actually used,
\begin{equation}\label{eq:pc-post-remainder-jets}
 \begin{aligned}
 \sum_{i=0}^1\sum_{j=0}^2
  |\mathcal E^i\partial_\lambda^j\mathcal R_1|
  &\le C q_1, & q_1&=Z^{-1}+z+m\varepsilon^3Z^{-3},\\
 \sum_{i=0}^1\sum_{j=0}^2
  |\mathcal E^i\partial_\lambda^j\mathcal R_2|
  &\le C q_2, & q_2&=\varepsilon^2(Z^{-1}+z)
                         +m\varepsilon^3Z^{-3}.
 \end{aligned}
\end{equation}
These bounds do not expand in $\varepsilon^2\log Z$. Increasing $Q$
and decreasing the fixed outer section makes the full field inward
pointing at the two fixed band boundaries and gives
\[
 f_\varepsilon'+\partial_\lambda\mathcal R_1\le-c_0<0
\]
throughout the interval.

At fixed $\tau$ and fixed field, differentiation in
$\ell_f=\log z_f$ acts on the coefficients as $\mathcal E$. Thus
$Z_1=\partial_{\ell_f}\lambda(\tau;z_f)$ solves
\[
 (Z_1)_\tau=(f_\varepsilon'+\partial_\lambda\mathcal R_1)Z_1
                    +\mathcal E\mathcal R_1.
\]
Its initial value is $O(m\varepsilon^2)$. With $\tau_Q$ defined and
truncated as above, the preceding transition interval has bounded length and bounded
variational coefficients, so
$\sup_{[0,\tau_Q]}|Z_1|+\int_0^{\tau_Q}|Z_1|\,d\tau\le C$.
On the remaining interval, when nonempty, variation of constants
gives explicitly
\[
 |Z_1(\tau)|\le Ce^{-c_0(\tau-\tau_Q)}
   +C\int_{\tau_Q}^{\tau}e^{-c_0(\tau-s)}q_1(s)\,ds.
\]
Moreover, with $\tau_o=\log(\xi_{\rm out}/z_f)$,
\[
 \int_0^{\tau_o}q_1\,d\tau
 \le \int_{Z_f}^{\infty}\frac{dZ}{Z^2}
      +\int_{z_f}^{\xi_{\rm out}}dz
      +m\varepsilon^3\int_{Z_f}^{\infty}\frac{dZ}{Z^4}\le C,
\]
because $Z_f\ge c_r$.  Fubini's theorem in the preceding stable-kernel
formula now yields
\begin{equation}\label{eq:pc-post-variation-bound}
 \sup_{[0,\tau_o]}|Z_1|+\int_0^{\tau_o}|Z_1|\,d\tau\le C.
\end{equation}

Put $\psi=z\mathcal C$. The full post-hit derivative, including
the moving terminal point, is
\[
 \partial_{\ell_f}\int_0^{\tau_o}\psi\,d\tau
 =-\psi(\tau_o)+\int_0^{\tau_o}
              (\psi_\lambda Z_1+\mathcal E\psi)d\tau.
\]
On the stable interval,
$\psi_\lambda=-2\varepsilon^2(\lambda+\varepsilon^2)^{-3}
+\partial_\lambda\mathcal R_2$ and
$\mathcal E\psi=\mathcal E\mathcal R_2$.
Equations \eqref{eq:pc-post-remainder-jets} and
\eqref{eq:pc-post-variation-bound} bound the integral; its truncated
pre-stable part has bounded length and smooth coefficients. A moving intermediate
endpoint appears with opposite signs in the two parts and cancels.
At the terminal section the normalized ordinate has derivative
$Z_1(\tau_o)-\lambda_\tau(\tau_o)$, and hence
\[
 \partial_{\ell_f}y(\xi_{\rm out};z_f)
 =P(\xi_{\rm out})[Z_1(\tau_o)-\lambda_\tau(\tau_o)].
\]

\medskip\noindent\textit{Completion on the returning interval.}\quad
Finally \eqref{eq:pc-energy-defect} gives the exact factorization
\[
 N(z_0)=(E_b)_y(z_0,y_0(z_0))
  \exp\left(\int_{z_0}^{\xi_{\rm out}}\mathcal C\,dz\right)
  U_{\rm reg}(y(\xi_{\rm out};z_0),p),
\]
where $U_{\rm reg}>0$ is the derivative of the remaining regular
holonomy in the $E_b$ coordinate.  The entrance factor has bounded
logarithmic $\log z_0$ derivative.  The pre-hit calculation bounds both
the defect derivative and $\partial_{\log z_0}\log z_f$; the post-hit
calculation above bounds the remaining defect derivative and terminal
ordinate derivative.  For the already fixed pair of physical sections,
$U_{\rm reg}$ is a positive $C^2$ unit after the final parameter
shrinking.  These bounds prove \eqref{eq:pc-normal-bound}.  Constants may
depend on the fixed sections; $L$ and the final parameter neighborhood
are chosen afterward, so no reciprocal section size is absorbed into
the small factor used below.

Every integral in this proof starts at the current point and proceeds
rightwards, except the boundary kernel, which also uses only its
right-hand tail. If $z\ge rZ_b$ lies beyond all roots, the comparisons
and Euler estimates used above hold on that tail. In particular
$p_a'\le2Z$ and the positive sum in \eqref{eq:pc-growing-ratio} are
unchanged. No margin to the left of $rZ_b$ is required. This proves
the second assertion.
\end{proof}

The curvature follows from a moving departure section, not from a
uniform bound for $N$ itself. Put $W=y-y_\sigma$, $K=sW_s/W$.
The divided coefficient relative to this boundary solution is
\[
 \mathfrak a(z,W)=\int_0^1H_y(z,y_\sigma+tW)\,dt,
 \qquad W_z=\mathfrak a(z,W)W.
\]
Here $W(0)=\kappa^2s$, so $K(0)=1$.
Before the first hit $W=\theta S$, exact differentiation gives
\[
 (\log K)_z=W\mathfrak a_W,\quad
 D=\mathfrak a-P'/P\asymp\frac z{\varepsilon^2P},\quad
 |W\mathfrak a_W|+|W^2\mathfrak a_{WW}|\le C(W/S)D
\]
on $z\ge L\kappa$. The bounded initial $\zeta$-interval has the
same integrated bound. In particular $K=1+O(\theta)$ before the hit.
Differentiating its integral equation gives the exact total derivative
\[
 s\partial_s\log K_g
 =s(z_g)_s[W\mathfrak a_W]_g+
   \int_0^{z_g}K(W\mathfrak a_W+W^2\mathfrak a_{WW})\,dz.
\]
The integral is $O(\theta)$ since $(W/S)_z=(W/S)D$, and the endpoint term
has the same bound by the differentiated hitting equation. Thus
\begin{equation}
\label{eq:pc-gate-identities}
 K_g=1+O(\theta),\quad |s(\log K_g)_s|\le C\theta,\quad
 (z_g)_s=-\frac{K_g}{sD_g},\quad
 M(s)=N(z_g)\theta S(z_g)\frac{K_g}{s}.
\end{equation}
The last identity uses the tangent to the first-intersection curve,
$y_g'-H(z,y_g)=-\theta SD$ and therefore includes the change of
hit time. Lemma~\ref{lem:pc-normal-derivative} now gives
\[
 -s(\log M)_s
 =1-s(\log K_g)_s+
     \frac{K_g}{D_g}\left((\log N)_z+\frac{P'}P\right)_{z_g}.
\]
Here $(z_gD_g)^{-1}\le
C(\varepsilon^2+\varepsilon/L+L^{-2})$. First choose $\theta$ small,
then $L$ large and $\varepsilon$ small. The last display is at least
$3/4$.
An inner hit $z_g\le L\kappa$ implies $s\ge c_L\theta$ by the
bounded inner difference equation, so it is absent when $s\le s_0$
for a sufficiently small fixed $s_0$.

To obtain the weighted bound, write $q=\sigma+\eta$, $J=q_s$,
$D_0=\mathsf B-\mathsf C(q+\sigma)>c$. Direct subtraction gives
\begin{equation}
\label{eq:pc-signed-cancellation}
 \frac d{d\zeta}\log\frac{sJV(\zeta,q)}{\eta V(\zeta,\sigma)}
 =\frac{\eta}{V(\zeta,q)}(\mathsf B'-\mathsf C\sigma')
 -\frac{D_0\eta}{V(\zeta,q)}(\log V(\zeta,\sigma))'.
\end{equation}
The first term has bounded integral; the coefficient of the second
belongs to $(0,1)$. Under \eqref{eq:pc-separated-polynomial} the
negative variation of $\log V(\zeta,\sigma)$ is bounded, because
$\mathsf A$ stays positive and $\sigma$ decreases. The initial
quotient is a positive bounded unit by $q_0\le Q_0$. At the outer
section,
\[
 sJ\le C V(\zeta,\sigma)/D_0\le C(1+\xi_{\rm out}^2/r^2),
 \qquad sM\le C\varepsilon^2(r^2+\xi_{\rm out}^2).
\]
The signed-cancellation identity therefore gives the size bound on the
whole actual returning interval $I$. Fix the incoming and short-map
constants and choose a tolerance $\delta=\delta_*>0$ satisfying both
inequalities in~\eqref{eq:pc-ratio-smallness}. Then choose the two sections
sufficiently small with their fixed ratio and shrink $r$ to make the size
coefficient below $\delta_*$. The short-map constants can be fixed in advance
as explained in Lemma~\ref{lem:pc-event-gluing}. After these sections
fix $L$ and $s_0$, use $w_\beta\le7m$ and the exponential wedge
\eqref{eq:pc-exponential-wedge} to shrink $\varepsilon$ until
\eqref{eq:pc-ratio-localization} holds, so that
$R\ge1/2$ implies $s\le C\delta w_\beta<s_0$.
Lemma~\ref{lem:pc-ratio} now applies with both sections fixed before $s_0$. It bounds the total isolated multiplicity by two in this
separated root-free case.

\subsection{Two separated positive roots}

For separated slow roots, we use an exact saddle passage and estimate
its complete complement. The comparison then counts every retained
return of a fixed field on the common physical source interval.

Suppose $-M\le a\le-2-\delta_a$. Put
\[
 A=|a|r,\quad c=a^{-2},\quad n=Ac,\quad
 e=\varepsilon/\sqrt c,\qquad
 X=z/n,\quad Y=y/(n^2e^2).
\]
Here $c$ lies in a compact subinterval of $(0,1/4)$, and
\begin{equation}
\label{eq:pc-two-root-system}
 X'=e^2(YQ+1-X+cX^2),\qquad Y'=XY-em,\qquad
 Q=1+bnX-n^2e^2Y.
\end{equation}
Write $p_c(X)=1-X+cX^2$. Its smaller reference root $X_s$ is a positive unit,
$p_c'(X_s)=-\sqrt{1-4c}$ is a negative unit, and the other root
$X_2$ is separated from it. The actual equilibria instead satisfy
\[
 Y=em/X,\qquad
 Q_e(X):=p_c(X)+\frac{em(1+bnX)}X
                      -\frac{n^2e^4m^2}{X^2}=0.
\]
On this compact $c$-interval the implicit function theorem gives
$X_i^{\rm eq}=X_i+O(em)$, with common finite phase bounds.
Their separation is at least half the reference gap after shrinking,
and $Q_e'(X_s^{\rm eq})$ remains a negative unit. All saddle-centred
coordinates and eigenvalues below refer to this actual equilibrium.
Fix a small collar width $\sigma$ below
one tenth of their minimum gap. Apply the graph construction in
Lemma~\ref{lem:small-ratio-saddle} on this fixed saddle collar.
The normal eigenvalue is a positive unit, all base derivatives contain
$e^2$, and $p_c'(X_s)$ stays away from zero. Thus
Lemma~\ref{lem:uniform-normal-product} and its simple-root scalar
change give, at fixed incoming and outgoing sections, the exact $C^j$
Dulac map
\begin{equation}
\label{eq:pc-two-root-power}
 D(s)=\xi_0(s/\omega_0)^{\rho_s},\qquad
 \rho_s=e^2\beta_s,\quad 0<c_s\le\beta_s\le C_s.
\end{equation}
The fixed spectral and root gaps give this adaptation of the
small-$c$ construction uniformly on the present compact $c$-interval.

The smallness of $c$ used near a simple contact must here be replaced
by smallness of $A$. For $x=-z\ge Ln$,
\[
 P(-x)=x^2+Ax+A^2c\asymp x(x+A),\qquad
 T_A=\log\frac{x_g+A}{A(1+Lc)}=\log(1/A)+O(1).
\]
On the upper barrier, the negative graph derivative satisfies
$H_y\le-\gamma/[\varepsilon^2(x+A)]$. Its integrated contraction
is at most $-\gamma T_A/(ce^2)$. By making $A$ small, this absorbs
any fixed loss $C/e^2$ from a bounded scaled interval. It absorbs
polynomial scaling losses too: for every fixed $j$,
\[
 A^{-j}\left(\frac{A(1+Lc)}{x_g+A}\right)^{\gamma/(ce^2)}
\]
is increasing on $0<A\le A_0$ once $e$ is sufficiently small.
There is no lower bound on $A$.
Since $d<0$, the incoming scaled height is consequently less than
$e^{-B/e^2}$ for any fixed $B$, after $A_0$ is chosen. It decreases
on the negative half and grows by at most $e^{C_\sigma/e^2}$ up to
$X_s-\sigma$. Thus every positive returning trajectory reaches the
saddle section below its fixed positive normal outlet; an earlier
hit of $Y=0$ is terminal. There is no early bypass in this case.
Only $s>0$ in \eqref{eq:pc-two-root-power} reaches the positive
outlet: $s=0$ lies on the stable separatrix, whereas for $s<0$ the
expanding negative normal deviation crosses $Y=0$ while the
separatrix height remains $O(em)$.

The outgoing trajectory must still pass the second-root region. From
the positive saddle outlet use $Y$ as phase:
\[
 \frac{dX}{dY}=\frac{e^2(YQ+p_c(X))}{XY-em},\qquad
 \frac{d(X^2)}{dY}\le Ce^2(1+1/Y)+Cce^2X^2/Y.
\]
These inequalities bound $X$ up to $Y=e^{-2}$. Before a fixed large
$Y_0$, its change is $O(e^2(1+|\log Y_0|))$, so it stays to the
left of $X_2$. Choose $Y_0$ above four times the bound for $X$;
then $YQ+p_c\ge Y/4$. A first zero of this speed is impossible:
at such a zero its $Y$-derivative is $1+bnX-2n^2e^2Y>0$.
Thus the orbit reaches $y=n^2$ with $z/n$ in a fixed positive
interval. The invariant-interval argument for $y/z^2$ in
\eqref{eq:pc-cone-equation} below gives its positive parabolic
continuation, independently of either polynomial root.

Let $C_+$ be this whole outgoing map, and $N_-$ the remaining fast
arc, decreasing-height entrance, negative passage and incoming section
conversion. The proof of Lemmas~\ref{lem:positive-transition} and
\ref{lem:negative-transition}, with the replacements just given,
yields
\[
 C_+'>0,\quad C_+'\le Cn^{3/2},\quad |C_+''/C_+'|\le C,\qquad
 N_-'>0,\quad |N_-''/N_-'|\le C\varepsilon^{-2}.
\]
Indeed the positive $Y$-phase derivatives through order three have
majorant $Ce^2(1+1/Y)$; the later relative-normal equation is unchanged.
On the negative side the first and second variations have a polynomial
scaling loss times $e^{-\gamma T_A/(ce^2)}$, followed by at most
$e^{C/e^2}$. The displayed negative action absorbs both. Hence, for
the full complement $H=N_-\circ C_+$,
\[
 \rho_s|H''/H'|\le Ce^2+CA^{3/2}\sqrt c=o(1).
\]
Lemma~\ref{lem:small-ratio-saddle} now gives total isolated
multiplicity at most two for $D\circ H$ on its complete return interval.

\subsection{The positive double root}

Near the positive double root, we need coordinates that remain uniform
as the slow roots coalesce. We first select the boundary graph and
control its phase derivatives, then straighten its normal flow. The
resulting product retains the full passage multiplier for the weighted
comparison.

It remains to treat $a$ in a fixed small neighborhood of $-2$.
In $z=rZ$, $y=r^2\varepsilon^2Y$, put $w=Y-\varepsilon m/Z$ and
rescale time by the positive factor $r$; this changes neither trajectories
nor return maps.  In this rescaled time variable,
\begin{equation}
\label{eq:pc-collision-system}
 \dot Z=\varepsilon^2F_s(Z,w),\qquad
 \dot w=Zw+\frac{\varepsilon^3m}{Z^2}F_s(Z,w),\qquad
 F_s=Q_\varepsilon+A_\varepsilon w-r^2\varepsilon^2w^2,
\end{equation}
where
\[
 Q_\varepsilon=p_a(Z)+\frac{\varepsilon m(1+brZ)}Z
               -\frac{r^2\varepsilon^4m^2}{Z^2},\qquad
 A_\varepsilon=1+brZ-\frac{2r^2\varepsilon^3m}{Z}.
\]
Choose fixed bases $Z_-<1-2\delta<1+2\delta<Z_+$ in the positive
half-plane, with fixed buffers on either side. All possible roots
of $Q_\varepsilon$ in this strip lie strictly between the bases.
Write $Z_b=Z_+$ for the right exterior base below.
Lemma~\ref{lem:uniform-normal-product} requires a chosen invariant
graph with finite phase bounds. On the no-root side we must obtain those
bounds for the solution selected by the physical boundary condition,
even when the slow speed has an arbitrarily small minimum.

\begin{lemma}[No-root tail and collision continuation]
\label{lem:pc-no-root-tail}
Fix a finite phase order $k$ (at least $2j+11$ for the subsequent
full normal-product output of order $j+2$) and choose fixed numbers

\[
 0<Z_-<Z_b<Z_1<Z_2
\]

and a fixed $\delta_0>0$ with $Z_\ell:=Z_--\delta_0>0$, so that
$Q_\varepsilon$ is a positive unit on $[Z_\ell,Z_-]$, every possible
root in the collision strip lies in $(Z_-,Z_b)$, and
$P(rZ)/(r^2Z^2)$ is a positive unit for $Z\ge Z_b$.  There are constants
$q_*>0$, $K_k>0$, and one parameter shrinking, independent of the minimum
of $Q_\varepsilon$, such that on the no-root side

\[
 Q_\varepsilon>0\quad\hbox{on }[Z_\ell,Z_2],\qquad
 Q_\varepsilon\ge q_*\quad\hbox{on }[Z_b,Z_2],
\]

the physical boundary solution $y_\sigma$ from
\eqref{eq:pc-boundary-tail} induces a unique graph $w=h(Z)$ on
$[Z_\ell,Z_1]$ and

\[
 -K_k\varepsilon^3m\le h(Z)<0,\qquad
 \max_{0\le i\le k}\|\partial_Z^ih\|_{L^\infty[Z_\ell,Z_1]}
 \le K_k\varepsilon^3m .
\]

At $m=0$ the graph is $h\equiv0$.  No constant in this statement uses a
positive lower bound for $Q_\varepsilon$ on $[Z_-,Z_b]$.
\end{lemma}

\begin{proof}
The case $m=0$ is immediate, so assume $m>0$.
Put $\tau=\varepsilon^2$, $n=\varepsilon^3m$, and
$\nu=r^3n$.  On the physical right tail $z\ge rZ_b$ consider
\[
 U(z)=\frac{\nu}{z},\qquad
 L(z)=(1-M\varepsilon^2)\frac{\nu}{z},
\]
where $M$ is fixed large before $\varepsilon$ is decreased.  Since
$P(rZ)/(r^2Z^2)$ is a positive unit there, the strip $L\le y\le U$
satisfies, uniformly in $r/\varepsilon$,
\[
 c\varepsilon^2z^2\le
 y(1+bz-y)+\varepsilon^2P(z)\le C\varepsilon^2z^2.
\]
For the scalar equation $y_z=H:=(zy-\nu)/F_{\rm ph}$, the reverse-flow
boundary expression $-H+J'$ is negative at $J=U$ and positive at
$J=L$ once $M>2C$.  Thus the strip is reverse invariant.  Terminal
solutions at $z=T$ with
$y(T)=\nu/((1+\varepsilon^2)T)$ remain in the strip; on it
\[
 (H)_y=
 \frac{z(y^2+\varepsilon^2P)+\nu(1+bz-2y)}{F_{\rm ph}^2}
 \ge\frac{c}{\varepsilon^2z}.
\]
The scalar mean-value equation therefore gives, for $T'>T$ and fixed
$z<T$,
\[
 |y_{T'}(z)-y_T(z)|\le \frac{\nu}{T}(z/T)^{c/\varepsilon^2}.
\]
The terminal solutions converge on every compact tail interval to a unique
positive solution. To identify its boundary condition, set $\ell=\log z$
and $v=zy/\nu$. The strip bounds $v$, and, with the field fixed,
$F_{\rm ph}/z^2=\varepsilon^2+O(e^{-\ell})$ gives
\[
 v_\ell=(1+\varepsilon^{-2})v-\varepsilon^{-2}+O(e^{-\ell}).
\]
Variation of constants for the bounded solution yields
$v=(1+\varepsilon^2)^{-1}+O(e^{-\ell})$, so this limit is
$y_\sigma$ in~\eqref{eq:pc-boundary-tail}. On the fixed right buffer
$F_{\rm ph}\ge c\varepsilon^2z^2$, so for each fixed finite $k$ the
differentiated scalar variational equations are regular and give $C^k$
convergence with the finite-order majorants recorded below.  Hence the
terminal jets of $y_\sigma$ exist and graph differentiation below is
justified.  Its pullback to the
collision coordinates satisfies
\[
h(Z)=\frac{y_\sigma(rZ)}{r^2\varepsilon^2}-\frac{\varepsilon m}{Z},
 \qquad -K n/Z\le h(Z)<0
\]
on the right buffer.  Enlarge $K$ by the fixed factor $1/Z_\ell$ if
necessary, so that this bound also implies $h\ge-Kn$ on the collision
rectangle below.

It remains to continue the graph through the possibly small slow speed.
In reverse time $s=-t$, the horizontal boundaries of the strip
\[
 Z_\ell\le Z\le Z_2,\qquad -Kn\le w\le0
\]
point inward while $F_s>0$, and the base coordinate decreases towards
$Z_\ell$.  Shrink the fixed tube once so that
$F_w\ge c_*>0$ there.  At $w=0$ one has
$w_s=-nF_s/Z^2<0$, while at $w=-Kn$ one has
$w_s\ge n(KZ_\ell-C/Z_\ell^2)>0$. On $F_s=0$ with $w<0$,
\[
 (F_s)_s=-F_wZw>0.
\]
The corner $F_s=w=0$ is impossible because there
$F_s=Q_\varepsilon>0$.  Thus the tail remains in the rectangle and
$Z_s=-\tau F_s<0$. If it did not reach $Z_\ell$ in finite reverse time,
compactness and monotonicity would give a nonempty invariant limit set on
$F_s=0$; the equation for $w$ then forces $w=0$, contradicting
$F_s=Q_\varepsilon>0$.  This proves continuation as a graph without any
bound on the crossing time.

On the right buffer, where $Q_\varepsilon\ge q_*$ and $|h|\le K n/Z_b$,
shrink the parameter box once more so that
\[
 |A_\varepsilon h|+r^2\varepsilon^2h^2\le q_*/2.
\]
Consequently $F_s\ge q_*/2$ throughout this buffer.  The graph equation
\[
 \tau F_s h_Z=Zh+\frac{n}{Z^2}F_s
\]
then gives the crude terminal bounds
\[
 |h^{(i)}(Z)|\le C_i n\tau^{-i}
       \qquad(Z\in[Z_1,Z_2],\ 0\le i\le k):
\]
the case $i=1$ follows directly, and differentiating the equation
successively gives the induction step because $F_s\ge q_*/2$ and all
coefficients on the fixed buffer are bounded.  Since $F_s$ is also
uniformly bounded above there, a reverse traversal of any fixed subinterval
of the right buffer takes at least a fixed positive multiple of
$\tau^{-1}$.  Fix the required finite order $k\ge1$ and partition this
buffer by
\[
 \zeta_j=Z_2-\frac{j}{k}(Z_2-Z_1),\qquad 0\le j\le k.
\]
Every subbuffer $[\zeta_j,\zeta_{j-1}]$ has reverse crossing time at
least $c_k/\tau$.  For $p=h_Z$, differentiation and
conversion to reverse time give
\[
 p_s=-\left(Z+\frac{nF_w}{Z^2}-\tau F_Z-\tau F_wp\right)p
      -h-n\left(\frac{F_Z}{Z^2}-\frac{2F}{Z^3}\right).
\]
On the first subbuffer the crude bound gives $\tau F_wp=O(n)$.
Thus the coefficient in parentheses is $Z+O(n+\tau)\ge Z_\ell/2$
after shrinking. Variation of constants across this subbuffer, whose
crossing time is at least $c_k/\tau$, gives $|p(\zeta_1)|\le K_1n$.
The interval $|p|\le K_1n$ is then inward pointing in reverse time,
so the good first-jet bound holds on $[Z_\ell,\zeta_1]$.

For $i\ge2$, differentiate the graph equation $i$ times and collect the
highest derivative.  With $p_i=h^{(i)}$ one obtains
\[
 \tau F_s(p_i)_Z=L_i p_i+R_i,\qquad
 L_i=Z+\frac{nF_w}{Z^2}
       -\tau\{iF_Z+(i+1)F_wh_Z\}.
\]
Every term in $R_i$ contains $n$, $h$, or a lower jet.
Inductively, suppose the jets through order $i-1$ have their good bounds
on $[Z_\ell,\zeta_{i-1}]$. Then $|R_i|\le C_i n$ throughout the next
subbuffer $[\zeta_i,\zeta_{i-1}]$ and $L_i\ge Z_\ell/2$.
At its right endpoint the crude estimate is still
$|p_i|\le C_i n\tau^{-i}$.  Variation of constants in reverse time
over this subbuffer gives
\[
 |p_i(\zeta_i)|\le C_i n\tau^{-i}
       e^{-c_kZ_\ell/(2\tau)}+C_i'n\le K_i n.
\]
The reverse barriers $|p_i|\le K_i n$ propagate this bound on
$[Z_\ell,\zeta_i]$. Thus each jet order uses its own subbuffer after
the lower-jet estimates have become available there.  After finitely
many steps all jets through order $k$ have the claimed bounds on
$[Z_\ell,Z_1]$.
This proves the stated $C^k$ estimate without dividing by the vanishing
minimum of $Q_\varepsilon$.
\end{proof}

On the rooted side, normal hyperbolicity supplies a graph across the
strip. For $m>0$, dividing the normal variable by $\varepsilon^3m$
retains the forcing factor in its phase bounds; the incoming separatrix
selects the
part needed by the return.

\begin{lemma}\label{lem:pc-selected-graph}
On the no-root side, the tail solution
\eqref{eq:pc-boundary-tail} continues through the fixed extended strip
$[Z_\ell,Z_1]$ as a graph $w=h(Z)$. On the rooted side there is such an
invariant graph whose incoming restriction is the stable separatrix
of the first simple root, or the incoming parabolic separatrix at
a double root. For every prescribed finite $k$, these graphs satisfy
$\|h\|_{C^k}\le K_k\varepsilon^3m$, uniformly as the roots coalesce.
The bounds concern phase derivatives; joint parameter differentiability
across the discriminant is not asserted.
\end{lemma}

\begin{proof}
Put
\[
 \tau=\varepsilon^2,\qquad \alpha=\varepsilon m,\qquad
 n=\varepsilon^3m=\tau\alpha,\qquad
 \varepsilon^4m^2=\tau\alpha^2.
\]
The last identity is needed in the coefficient of $Q_\varepsilon$.
Write also $\nu_s:=n$ and abbreviate $F:=F_s$.
For $n>0$ write $w=n\psi$.  The exact auxiliary system is
\begin{equation}\label{eq:pc-graph-auxiliary}
 \dot Z=\tau F(Z,n\psi),\qquad
 \dot\psi=Z\psi+Z^{-2}F(Z,n\psi),
\end{equation}
where
\[
 Q_\varepsilon=p_a(Z)+\frac{\alpha(1+brZ)}Z
                   -\frac{r^2\tau\alpha^2}{Z^2},\qquad
 A_\varepsilon=1+brZ-\frac{2r^2\tau\alpha}{Z}.
\]
Because $n=\tau\alpha$, \eqref{eq:pc-graph-auxiliary} extends smoothly
to $n=0$.  Freeze $(\alpha,a,b,r)$ in a compact box and let $\tau$ tend
to zero.  At $\tau=0$ its critical graph is
\[
 \psi_0(Z)=-\frac{Q_0(Z)}{Z^3},\qquad
 Q_0(Z)=p_a(Z)+\frac{\alpha(1+brZ)}Z.
\]
With $u=\psi-\psi_0$, the equations on the working strip are exactly
\begin{equation}\label{eq:pc-graph-auxiliary-shift}
 \dot Z=\tau F_\tau,\qquad
 \dot u=Zu+Z^{-2}(F_\tau-Q_0)-\tau\psi_0'F_\tau,\qquad
 F_\tau=F(Z,n(\psi_0+u)).
\end{equation}
Thus their $C^N$ distance from $(\dot Z,\dot u)=(0,Zu)$ is
$O_N(\tau)$ on one fixed tube, with no inverse power of $n$ and no
root-separation constant.

Choose a fixed cutoff equal to one on a larger collision strip, extend
$\psi_0$ by a fixed extension operator, and close the base into a circle.
Choose a positive function $a_0$, equal to $Z$ on the working strip and bounded below by a positive constant on
the circle, and extend \eqref{eq:pc-graph-auxiliary-shift} there as
\[
 \dot\theta=\chi\tau F_\tau,\qquad
 \dot u=a_0u+\chi\{Z^{-2}(F_\tau-Q_0)-\tau\psi_0'F_\tau\}.
\]
At $\tau=0$ the invariant circle $u=0$ has zero tangential rate and a
fixed normal gap. Fix the desired phase order $k$ first and choose
the source order at least $k+1$. Fenichel's
finite-order persistence theorem
\cite[Theorem~3, p.~215]{Fenichel1971} gives an invariant
$C^{k+1}$ graph $u=\phi$ close to zero. The common phase bounds are
obtained from its invariance equation. Write the extended field as
$\dot\theta=\mathcal A(\theta,u)$ and
$\dot u=\mathcal B(\theta,u)$; then
$\mathcal A=O_{C^{k+1}}(\tau)$ and
$\mathcal B=a_0u+O_{C^{k+1}}(\tau)$.
At a maximum of $|\phi|$, the equation
$\mathcal A(\theta,\phi)\phi'=\mathcal B(\theta,\phi)$ gives
$\|\phi\|_{C^0}\le C\tau$. Differentiate successively through
order $k$ and evaluate at a maximum of each derivative's absolute
value. The next derivative vanishes there; the coefficient of the
highest retained derivative is
$\mathcal B_u+O_k(\tau)\ge\min a_0/2$ after shrinking. The first
step uses the close graph's bounded first derivative, and all later
steps use the preceding bounds. The remaining terms are $O_k(\tau)$,
so induction gives $\|\phi\|_{C^k}\le C_k\tau$. Hence
$\psi=\psi_0+\phi$ has a common $C^k$ bound after one common
shrinking of $\tau$. Reconstructing $h=n\psi$ yields
\begin{equation}\label{eq:pc-graph-sharp-size}
 \|h\|_{C^k}\le K_kn=K_k\varepsilon^3m.
\end{equation}
For $n=0$ take $h=0$, which solves the original equation exactly.
The circle, the finite order, and $K_k$ are chosen before the parameter
shrinking, so \eqref{eq:pc-graph-sharp-size} remains uniform when
$\varepsilon^3m$ is smaller than every exponential in
$1/\varepsilon$.  Restricting the graph back to the original strip
removes all dependence on the auxiliary extension from the equations
used below.

For any such graph, invariance gives the exact factorization
\begin{equation}\label{eq:pc-graph-unit}
 \begin{gathered}
 Zh=(\varepsilon^2h_Z-\nu_s/Z^2)F(Z,h),\qquad
 F(Z,h)=UQ_\varepsilon,\\
 U=\left(1-\frac{(A_\varepsilon-r^2\varepsilon^2h)
                     (\varepsilon^2h_Z-\nu_s/Z^2)}Z\right)^{-1}>0.
 \end{gathered}
\end{equation}
In particular it contains the equilibria at the roots of
$Q_\varepsilon$. Its incoming orbit approaches the first root in
positive time. Every nonzero displacement from the graph grows at a
uniform positive normal rate while it stays in the tube. Hence this
incoming orbit is precisely the stable separatrix at a simple first
root, and the incoming parabolic separatrix at a double root. This
restriction is independent of the auxiliary extension. The rest of
the graph is used only to define coordinates on the strip.

On the no-root side, apply Lemma~\ref{lem:pc-no-root-tail} with the
same right buffer and collision bases. It continues the physical tail
\eqref{eq:pc-boundary-tail} with its sign and prescribed phase jets,
without dividing by the minimum of $Q_\varepsilon$. The factorization
\eqref{eq:pc-graph-unit} therefore holds for this graph, which is the
one used in the through passage. On the rooted side, the incoming
stable or parabolic separatrix is intrinsic; the remaining weak graph
may depend on the circle extension and serves only to define product
coordinates. It need not agree with a separately prescribed outgoing
tail. The rooted fibre-exit estimates use its finite phase jets and the
same incoming threshold.
\end{proof}

On an incoming branch the selected graph has the exact formula
\begin{equation}
\label{eq:pc-canonical-graph}
 h(Z)=-\varepsilon m\int_Z^{Z_T}
 \exp\left(-\int_Z^t\frac{x}{\varepsilon^2F_s(x)}\,dx\right)
                                  \frac{dt}{t^2},
\end{equation}
where $Z_T$ is the first root or infinity and $F_s(x)$ is evaluated on
the selected graph. The kernel is strictly between zero and one, so
$-\varepsilon m/Z<h<0$ there.
The positive trace $Y_\sigma=h+\varepsilon m/Z$ pulls back uniquely
to $q_0=\beta>0$. On the no-root side it is exactly the same
trajectory as \eqref{eq:pc-boundary-tail}, not another choice of slow
manifold. Hence the two definitions of $\beta$ agree on their overlap.
On the outgoing side of this threshold, the physical displacement is
$W=y-y_\sigma>0$, with $W(0;s)=\kappa^2s$. On the collision strip,
this is $W=\kappa^2(w-h)$, not the nullcline offset $w$.

Translate this graph by $u=w-h(Z)$. The base equation is
$\dot Z=\varepsilon^2v(Z)+\varepsilon^2c(Z,u)u$, and the divided
normal coefficient is
\[
 B(Z,u)=Z+\left(\frac{\varepsilon^3m}{Z^2}
                    -\varepsilon^2h_Z\right)c(Z,u)\ge c_*>0.
\]
Fix nested compact intervals
\[
 [Z_-,Z_b]\Subset I_0\Subset I_1\Subset I_2\Subset(Z_\ell,Z_1).
\]
They and the physical bases $Z_-,Z_b$ are fixed before the finite-order
parameter shrinking. The selected graph and its phase jets are already
defined on the larger interval $[Z_\ell,Z_1]$, so the product is taken
over $I_0$ with both physical bases in its interior.
Lemma~\ref{lem:pc-selected-graph} supplies the finite jet bounds
required in Lemma~\ref{lem:uniform-normal-product}, with
$\tau=\rho=\varepsilon^2$. Fix the full output order $J=j+2$ first;
the selected graph is constructed once with $k\ge2j+11$, leaving
$C^{2j+10}$ coefficients after $h_Z$. The product is then taken
at order $J$, with its section losses already reserved. The lower
order $2j+9$ sufficient for the fibre comparison alone is not used
for this full normal construction. Its exact coordinates give
\begin{equation}
\label{eq:pc-collision-product}
 \dot\theta=\varepsilon^2v(\theta),\qquad
 \dot\eta=b_n(\theta)\eta,\qquad b_n\ge c>0.
\end{equation}
The normal coordinate is tangent to the identity and both coordinate
maps have uniformly bounded finite derivatives. At either exterior
base $v$ is a positive unit. The first section consequence of that
lemma gives conversion to a fixed $Z$-section in time $O(\eta)$,
with bounded derivatives, after cancelling the common
$\varepsilon^2$ in the hitting equation. No such cancellation at a
base root is needed.

To apply one return comparison, we must first place every counted cycle
in the same small source interval before the collision. The next lemma
obtains this from contraction along the preceding negative segment.

\begin{lemma}[Localization of the collision source]
\label{lem:pc-collision-localization}
Fix the collision coefficient interval, the exterior bases with their
buffers, and a forcing bound $m_0>0$, with $0<m\le m_0$. For every fixed
$B>0$ there are $r_0,\varepsilon_0>0$ such that every counted cycle
satisfies
\begin{equation}
\label{eq:pc-collision-localization}
 0<Y(Z_-)\le e^{-B/\varepsilon^2}
 \qquad(0<r<r_0,\ 0<\varepsilon<\varepsilon_0).
\end{equation}
For a fixed sufficiently small fiber height $\eta_0$, all these cycles
belong to one open interval $I_{\rm ret}$ of actual physical returning
sources on which the canonical left-base amplitude $A$ is defined and
$0<A<\eta_0$. They lie in its smaller part $A<\eta_0/2$.
\end{lemma}
\begin{proof}
Choose a fixed large $L$ and keep the left exterior base and its buffer
strictly to the left of the collision roots. After decreasing the fixed
coefficient interval, $p_a\ge p_*>0$ on $[-L,Z_-]$. The constants in
the negative-entry estimate of Lemma~\ref{lem:pc-entry} are common on
this coefficient box and the fixed range $0<m\le m_0$; thus
\[
 y(-Lr)\le K\varepsilon^3(Lr/x_g)^{\gamma/\varepsilon^2}.
\]
Here $K,\gamma,x_g>0$ and all sections are fixed before $r$ and
$\varepsilon$ are decreased. In the physical scaled variables the exact
orbit equation is
\[
 Y_Z=\frac{ZY-\varepsilon m}
 {\varepsilon^2\{p_a(Z)+(1+brZ-r^2\varepsilon^2Y)Y\}}.
\]
In a provisional fixed tube $0<Y<h_*$ the denominator in braces is
at least $p_*/2$. The height decreases for $Z\le0$, and for
$0\le Z\le Z_-$ it satisfies $Y_Z\le CY/\varepsilon^2$. Consequently,
before any first exit from the tube,
\[
 Y(Z)\le K\varepsilon r^{-2}
 \exp\left(-\frac{\gamma\log(x_g/(Lr))-C_+}{\varepsilon^2}\right)
 \qquad(-L\le Z\le Z_-),
\]
where $C_+$ is fixed. First choose $r_0<x_g/L$ so that
$\gamma\log(x_g/(Lr_0))>C_++B+1$. For
$N=\gamma/\varepsilon^2>2$, the function
$r^{-2}(Lr/x_g)^N$ increases with $r$. Hence the displayed bound is at
most
\[
 K\varepsilon r_0^{-2}e^{-(B+1)/\varepsilon^2}
 \le e^{-B/\varepsilon^2}
\]
after decreasing $\varepsilon_0$. Decrease it also so that this bound
is less than $h_*/2$. This puts the entrance inside the tube and excludes
a first upward exit. A counted cycle cannot hit $Y=0$, by
Lemma~\ref{lem:sections}. Since $\dot Z\ge\varepsilon^2p_*/2$ in the
rescaled time, it reaches $Z_-$ in finite physical time at most
$2(L+Z_-)/(r\varepsilon^2p_*)$. This proves
\eqref{eq:pc-collision-localization} without a relation between $r$ and
$\varepsilon$.

We now restrict the actual source interval. Let $I_{\rm phys}$ be the
positive contact-passage interval supplied by Lemma~\ref{lem:sections},
with the prescribed upper source core and common outgoing target core.
Every such passage has $F>0$, so its physical height $Y_-(s)$ at
$Z=Z_-$ is defined. The scalar variational equation gives
$\partial_sY_->0$ for the incoming height coordinate $s=q_0-\beta$
and each fixed field. The selected graph has physical trace
$Y_\sigma(Z_-)>0$, and its pullback to $z=0$ is $q_0=\beta$.
Thus $Y_--Y_\sigma(Z_-)>0$ is equivalent to $s>0$. On the no-root
side the selected tail stays below the outgoing fast core, and scalar
order excludes sources at or below it. On the root side the selected
incoming graph tends to the first equilibrium. The absolute value of a
negative normal difference grows at a positive rate by
\eqref{eq:pc-collision-product}, whereas the graph height is
$O(\varepsilon m)$; such a positive passage
hits $Y=0$ before leaving the fixed normal tube. The equality source
remains on the selected graph. Hence a counted cycle must have $s>0$.

Use the same selected graph, exterior base, and buffers as in
\eqref{eq:pc-collision-product}. The change from the physical section
$Z=Z_-$ to the canonical section $\theta=\theta_-$ gives a function
$\phi_p$ of the normal difference $u=Y_--Y_\sigma(Z_-)$, with
\[
 A=\phi_p(u),\qquad \phi_p(0)=0,\qquad c\le\phi_p'\le C
\]
on one common interval about zero. Indeed, $u=\eta+\eta^2L_2$ there;
the phase shear is $O(\varepsilon^2\eta)$ and the exterior base speed
is $\varepsilon^2$ times a positive unit. The first section consequence
of Lemma~\ref{lem:uniform-normal-product} gives an $O(\eta)$ flight
with bounded derivatives. It preserves the selected graph, so its normal
multiplier is a positive unit. Choose $\eta_0$ small enough to lie in
the common positive image of $\phi_p$, and put
$u_*(p)=\phi_p^{-1}(\eta_0)$. Define
\begin{equation}
\label{eq:pc-localized-source}
 I_{\rm ret}=\{s\in I_{\rm phys}:
       0<Y_-(s)-Y_\sigma(Z_-)<u_*(p)\}.
\end{equation}
This is an open interval, possibly empty, because $Y_-$ is strictly
increasing; the target restriction is already included in $I_{\rm phys}$.
On it $0<A<\eta_0$ and $A_s>0$. Every counted cycle is on the outgoing
side of the selected boundary and satisfies
\[
 0<A\le C(Y_--Y_\sigma(Z_-))\le CY_-
       \le Ce^{-B/\varepsilon^2}<\eta_0/2
\]
after one further decrease of $\varepsilon_0$. Thus all its source
values lie inside $I_{\rm ret}$. The artificial boundary $A=\eta_0$
contains no counted zero. The boundary $A=0$ is the selected invariant
trajectory and does not reach the outgoing fast core.
\end{proof}

On this actual source interval the difference equation before the left
base, where $p_a$ has a fixed positive margin, gives
\begin{equation}
\label{eq:pc-left-elasticity}
 K_A:=sA_s/A=1+O(\eta_0),\qquad
 |s(\log K_A)_s|\le C\eta_0.
\end{equation}
To see both bounds, use this positive displacement and its divided
coefficient $\mathfrak a$. Integrate
$(\log(sW_s/W))_z=W\mathfrak a_W$ and its $\log s$ derivative; their
majorants are $C\eta\,d\log\eta$, whose integral is $O(\eta_0)$.
The final amplitude coordinate is tangent to the identity, so it
preserves these estimates.

A fiber exit near a root still has to reach the regular outgoing arc.
We first follow the entire positive fiber to the height $y=r^2$, with
bounds uniform as the roots coalesce. Only then does the parabolic
coordinate $y/z^2$ enter a fixed contracting interval.

\begin{lemma}[Finite fibre-to-parabolic transition]
\label{lem:pc-finite-fibre-transition}
Fix a sufficiently small $\eta _0>0$ and a compact interval
$I_Z=[Z_-,Z_+]\Subset(0,\infty)$ containing the collision bases.  Let
$\eta=\eta _0$ denote the fibre in the exact product coordinates
\eqref{eq:pc-collision-product}.  Write $x=Z$ on its physical image and
$\Theta(x)$ for the corresponding product phase.  The physical fibre graph
is
\[
 w_0(x)=h(x)+\eta _0+\eta _0^2L_2(\Theta(x),\eta _0),
 \qquad
 \frac{\eta _0}{2}\le w_0(x)\le\frac{3\eta _0}{2},
\]
after decreasing $\eta _0$ once, with
$\|w_0-\eta _0\|_{C^3}\le C(\eta _0^2+\varepsilon^3m)$.
There are a coefficient neighbourhood $|a+2|<\delta_a$, a number
$r_0>0$, and $\varepsilon _0>0$ such that, for

\[
 |a+2|<\delta_a,\quad 0<r<r_0,\quad 0<\varepsilon<\varepsilon _0,
 \quad 0\le m\le m_0,
\]

every solution of the collision equations starting on this fibre with
$x=Z\in I_Z$ has a unique graph representation $Z=Z(Y)$ up to
$Y=\varepsilon^{-2}$.  On this interval

\[
 Y_{\rm in}(x)=w_0(x)+\frac{\varepsilon m}{x},\qquad
 c_0\le Z(Y)\le C_0,\qquad
 F_s(Z,Y-\varepsilon m/Z)\ge c_0\eta _0,
 \qquad ZY-\varepsilon m\ge c_0Y .
\]

The terminal graph $Z_t(x)$, parameterised by its initial fibre coordinate
$x=Z(Y_{\rm in}(x))$, satisfies

\[
 0<c_1\le (Z_t)_x\le C_1,
 \qquad |\partial_x^j Z_t|\le C_1\quad (j=2,3),
 \qquad c_1\le Z_t\le C_1,
\]

with constants independent of $r,\varepsilon,m$ and of the distance from
$a$ to the double-root value.
\end{lemma}

\begin{proof}
Write $Y$ for the phase and put
\[
 H=1+brZ-r^2\varepsilon^2Y,
 \qquad N=YH+p_a(Z),
 \qquad D=ZY-\varepsilon m.
\]
On the fibre, $w=Y-\varepsilon m/Z=w_0(Z)$.  Since

\[
 p_a(Z)=(Z-1)^2+(a+2)Z,
\]

we may choose $\delta_a$ and then $r_0$ so that, uniformly on the fixed
bootstrap strip,
$H\ge3/4$, $A_\varepsilon\ge3/4$, and
$F_y=1+brZ-2r^2\varepsilon^2Y\ge1/2$ on the whole phase strip.
Choose a fixed $B>Z_+$ with $B^2-Z_+^2>3$. The same identity allows
us to decrease $\delta_a$ once more so that
$p_a(Z)\ge-\eta _0/8$ on the enlarged bootstrap interval
$[Z_-/2,B]$.  Then the exact cancellation after
substitution of $w=Y-\varepsilon m/Z$ gives
\[
 F_s=N=YH+p_a(Z)\ge Y/2\qquad(Y\ge\eta _0/2).
\]
Choose $\varepsilon _0$ so that $\varepsilon m_0\le Z_-\eta _0/8$.
Then $D=Zw_0(Z)\ge Z\eta _0/2$ at the initial fibre.  In the region
$Z\in[Z_-/2,B]$, $Y\in[\eta _0/2,\varepsilon^{-2}]$, and
$D\ge ZY/2$, direct substitution in the collision equations gives
\[
 \frac{dZ}{dY}=f(Y,Z):=\varepsilon^2\frac{N}{D},
 \qquad \dot Z=\varepsilon^2F_s,
 \qquad \mathcal G:=D=ZY-\varepsilon m,
 \qquad G_{\rm original}=r^3\varepsilon^2\mathcal G.
\]
The preceding choices give $F_s=N\ge Y/2$ and
$D=ZY-\varepsilon m\ge ZY/2$ on this region, so $Z_Y>0$ and both
physical speeds are positive.

For $0\le k\le3$, the rational equation has
\[
 |\partial_Z^k f(Y,Z)|\le C_k\varepsilon^2(1+Y^{-1}),
\]
with constants depending only on the fixed strip, coefficient box and
$\eta_0$.
Moreover
\[
 (Z^2)_Y=2\varepsilon^2
 \frac{H+p_a(Z)/Y}{1-\varepsilon m/(ZY)}
 \le \frac{2\varepsilon^2(1+|b|rB)}{1-\chi_m}
       +\frac{2P_B\varepsilon^2}{(1-\chi_m)Y},
 \qquad
 \chi_m:=\frac{4\varepsilon m_0}{Z_-\eta _0}<\frac12,
\]
where $P_B=\sup_{|a+2|<\delta_a,\,Z\in[Z_-/2,B]}|p_a(Z)|$.
Choose a margin $\sigma>0$ with
$2+\sigma<B^2-Z_+^2$.  First decrease $r_0$ so that the contribution of
$|b|rB$ to the first numerator is at most $\sigma/4$, and then decrease
$\varepsilon_0$ so that $\chi_m$ is small enough and the logarithmic term
below is at most $\sigma/2$.  Before these final choices, integration up to
$Y=\varepsilon^{-2}$ gives
\[
 Z^2(Y)\le Z_+^2+
 \frac{2(1+|b|r_0B)}{1-\chi_m}
 +\frac{2P_B\varepsilon^2}{1-\chi_m}
   \log\frac{2\varepsilon^{-2}}{\eta_0}.
\]
Choose first $r_0$ and then $\varepsilon_0$ so that the last bound is
at most $Z_+^2+2+\sigma<B^2$. This closes the upper boundary;
$Z_Y>0$ closes the lower one, and $D\ge c_0Y$ follows from
$\varepsilon m_0\ll Z_-\eta_0$.
This proves existence through $Y=\varepsilon^{-2}$.

It remains to control the fibre dependence.  The initial graph is
$Y_{\rm in}(x)=w_0(x)+\varepsilon m/x$, so its first three derivatives
are uniformly bounded on $I_Z$ and
$Y_{\rm in}'=O(\eta _0^2+\varepsilon^3m+\varepsilon m)$.
At fixed $Y$, the first variation obeys
\[
 V_Y=f_ZV,\qquad
 V(Y_{\rm in}(x))=1-f(Y_{\rm in}(x),x)Y_{\rm in}'(x).
\]
The initial value lies in $[1/2,3/2]$, while
\[
 \int_{\eta _0/2}^{\varepsilon^{-2}}
 \varepsilon^2\left(1+\frac1Y\right)dY\le C.
\]
Thus $c_1\le V\le C_1$.  The second and third differentiated equations have the same linear
coefficient $f_Z$ and forcing $f_{ZZ}V^2$ and
$3f_{ZZ}VV^{(2)}+f_{ZZZ}V^3$, respectively. The displayed integrable
majorant and bounded initial jets therefore give
$|V^{(2)}|+|V^{(3)}|\le C_1$.
The terminal bounds follow.
\end{proof}

The terminal graph from the finite transition supplies the initial
conditions for parabolic contraction. The next lemma continues the
whole fiber to the target and controls its first two derivatives.

\begin{lemma}
\label{lem:pc-fast-cone}
After $\eta_0$ is fixed, the coefficient neighborhood of $a=-2$
can be chosen so that every point of the entire fixed fiber
$\eta=\eta_0$, $\theta_-\le\theta\le\theta_+$, reaches the same
regular outgoing target. In the coordinate $x=Z=z/r$ restricted to
this fiber, its target map $\Phi$ satisfies
\begin{equation}
\label{eq:pc-fiber-target}
 0<|\Phi_x|\le C_\xi r^\alpha,\qquad
 |\Phi_{xx}/\Phi_x|\le C,\qquad \alpha>0.
\end{equation}
The target section $\xi$ is fixed before shrinking $r,\varepsilon$.
\end{lemma}

\begin{proof}
Choose the coefficient interval and small parameters as in
Lemma~\ref{lem:pc-finite-fibre-transition}.
Lemma~\ref{lem:pc-finite-fibre-transition}, with fiber coordinate
$x=Z$, gives $F,G>0$, a fixed compact positive $Z$-interval, and three
terminal derivatives at $Y=\varepsilon^{-2}$, uniformly up to the
double root. Its phase equation is
\[
 Z_Y=\varepsilon^2\frac{Y(1+brZ-r^2\varepsilon^2Y)+p_a(Z)}
                              {ZY-\varepsilon m}.
\]
The terminal bounds are therefore
\[
 z_t=rZ_t(x),\quad y_t=r^2,\qquad
 c\le (Z_t)_x\le C,\quad
 |\partial_x^jZ_t|\le C\ (j=2,3),\quad c\le Z_t\le C.
\]

From there put $u_f=y/z^2$, $\ell=\log z$,
$\pi=P/z^2$, and $\delta_z=\nu/z^3$. The exact equation is
\begin{equation}
\label{eq:pc-cone-equation}
 (u_f)_\ell=\Psi(\ell,u_f)
   =\frac{u_f-\delta_z}
            {u_f(1+bz-u_fz^2)+\varepsilon^2\pi}-2u_f.
\end{equation}
The terminal height $Y=\varepsilon^{-2}$ gives $y_t=r^2$ and
$u_f(\ell_t)=1/Z_t^2$, uniformly separated from zero. This transition
from the original fiber height $y=O(r^2\varepsilon^2)$ puts the orbit
in the region where the parabolic contraction applies. For $z\ge cr$,
the first two Euler derivatives of $\pi$ are bounded; roots of $P$ occur only in
the additive term $\varepsilon^2\pi$. Choose a fixed interval
$I=[u_-,u_+]$ containing all $1/Z_t^2$, with
$0<u_-<1/2<u_+$. After first fixing $I$ and then choosing the target
$\xi$ small, the faces of $I$ point inward, since uniformly on $I$
\[
 \Psi=\frac1{1+bz-u_fz^2}+O(\varepsilon^2+\varepsilon^3)-2u_f.
\]
Writing $D=u_f(1+bz-u_fz^2)+\varepsilon^2\pi$ and retaining the
$u_f$-dependence of its denominator gives the exact identity
\[
 \Psi_{u_f}
 =\frac{\varepsilon^2\pi+\delta_z(1+bz-2u_fz^2)
       +u_f^2z^2}{D^2}-2.
\]
Here $1+bz-u_fz^2\ge c>0$ and $D\ge c>0$ on the forward invariant
collar. Since $\pi$ is bounded and $\delta_z=O(\varepsilon^3)$ there,
we may shrink $\xi$ and then the parameters so that
$\Psi_{u_f}\le-c<0$ uniformly. All derivatives needed through order
three are bounded. Thus the orbit reaches $z=\xi$ with $F\ge c y>0$.

If $J$ is the vertical variation of \eqref{eq:pc-cone-equation}, then
$0<J(\ell_\xi,\ell_t)\le(z_t/\xi)^\alpha$.
The moving initial point must also be included. Its transverse factor is
\[
 (u_t)_x-\Psi(\ell_t,u_t)(\ell_t)_x
 =-\frac{u_t-\delta_{z_t}}
              {u_t(1+bz_t-u_tz_t^2)+\varepsilon^2\pi(z_t)}
                                  \frac{(Z_t)_x}{Z_t}.
\]
It is a negative unit with bounded two derivatives. The formula for
the initial-time derivative of the flow, and
$J_\ell=\Psi_{u_f}J$, therefore give the first inequality in
\eqref{eq:pc-fiber-target}. Logarithmic differentiation of the
variational exponential has an endpoint term bounded by
$|\Psi_{u_f}(\ell_t,u_t)|$ and an integral with factor
$\Psi_{u_fu_f}J$; its absolute integral is bounded by the stable
kernel. The first logarithmic derivative of the transverse factor is
also bounded. This proves the second inequality. The final fixed
holonomy in energy coordinates preserves both estimates.
\end{proof}

On the interval which exits through $\eta=\eta_0$, write $x=Z=z/r$
at this hit. The change from the canonical
$\theta$ to $x$ on that fixed fiber is a $C^3$ diffeomorphism with
positive unit derivative. Write $\widetilde v$ for $v$ multiplied by
this derivative and $\widetilde b_n$ for $b_n$ composed with the inverse
change; both and their first derivatives remain bounded.
Let $x_-$ be the image of $\theta_-$
under this change. The backward incoming amplitude is
\[
 \mathcal D(x)=\eta_0\exp\left(-\int_{x_-}^x
             \frac{\widetilde b_n(t)}{\varepsilon^2\widetilde v(t)}\,dt\right).
\]
The base speed is positive along the fibre-exit branch before its first
root. Thus $\mathcal D$ decreases in the physical coordinate $x=Z$,
and differentiation of $A(s)=\mathcal D(x(s))$ gives
\[
 |sx_s|=\varepsilon^2|\widetilde v|K_A/\widetilde b_n,\qquad
 M=|\Phi_x|\varepsilon^2|\widetilde v|K_A/(\widetilde b_n s).
\]
By Lemma~\ref{lem:pc-fast-cone},
$sM\le C_\xi\varepsilon^2r^\alpha$ uniformly on the whole
fiber-exit returning interval. Moreover
\[
 \begin{aligned}
 -s(\log M)_s
 &=1-s(\log K_A)_s-sx_s\partial_x\log|\widetilde v\Phi_x/\widetilde b_n|\\
 &\ge1-C\eta_0-C\varepsilon^2\ge1/2.
 \end{aligned}
\]
There is no division by a root separation in this estimate:
$sx_s\,\widetilde v_x/\widetilde v=O(\varepsilon^2\widetilde v_x)$.
With the target fixed, shrink $r$ so that $C_\xi r^\alpha$ is below
the tolerance required by Lemma~\ref{lem:pc-ratio}. These are the size
and curvature bounds for the fiber-exit piece in the common comparison
below, uniformly up to the double root.

\subsection{Passage through the double-root region without a fiber exit}

We now treat the through passages across the collision strip and join them
to the fiber exits on the localized source interval of
Lemma~\ref{lem:pc-collision-localization}.

On the no-root side, sufficiently small amplitudes cross the whole
strip before reaching the fiber. Between the fixed bases, the exact
product~\eqref{eq:pc-collision-product} gives
\begin{equation}
\label{eq:pc-through-multiplier}
 B=m_b A(s),\qquad
 m_b=\exp\left(\int_{\theta_-}^{\theta_+}
                         \frac{b_n(t)}{\varepsilon^2v(t)}\,dt\right).
\end{equation}
The multiplier is finite for each such field, but has no uniform upper
bound. Passage from $\theta_-$ to $\theta_+$ occurs for $0<A<\eta_0/m_b$.
We keep \eqref{eq:pc-through-multiplier} in the complete map.
Figure~\ref{fig:collision-passages} places this interval beside the
fiber-exit interval. Their possible common boundary is included in the
same complete return map.



At the right base $z_b=rZ_b$ all possible roots are to the left and
$P\asymp r^2+z^2$, $P'>0$ for $z\ge z_b$.
The canonical graph is the same positive decreasing solution
$y_\sigma$ constructed from the right-hand boundary condition.
The change from the canonical base to $z=z_b$ is
\[
 t=\frac{y-y_\sigma}{\kappa^2}=\Psi_b(B),\quad
 \Psi_b(0)=0,\quad c\le\Psi_b'\le C,\quad
 |\Psi_b''/\Psi_b'|+|\Psi_b'''/\Psi_b'|\le C.
\]
The strong-fiber displacement is $O(\varepsilon^2B)$ and the exterior
base speed is $\varepsilon^2$ times a positive unit. Dividing the
implicit hit equation by $\varepsilon^2$ proves these bounds without
inverting the long passage.

Start \eqref{eq:pc-signed-cancellation} at $z_b$ with initial
difference $t$. Its initial quotient is a bounded positive unit.
On this right tail $\mathsf A=P/r^2$ increases and the total variation
of $\mathsf B\sigma-\mathsf C\sigma^2$ is bounded by $C\sigma(z_b)$.
Hence its negative logarithmic variation is bounded independently of
the earlier polynomial minimum. For the complete target derivative
$M_t$ this gives
\[
 tM_t\le C\varepsilon^2(r^2+\xi_{\rm out}^2).
\]
Take $\eta_0$ sufficiently small so that the initial ratio $W/S$ is
less than $\theta/2$. Retained cycles have that ratio larger than
$\theta$ at $\xi_{\rm in}$, so they first meet $W=\theta S$
before the inner section. Lemma~\ref{lem:pc-normal-derivative} applies
on this exterior tail, with room to complete the flight before
$\xi_{\rm out}$. In \eqref{eq:pc-gate-identities}, now with initial
variable $t$, the last logarithmic term is bounded by
$C/(z_gD_g)\le C\varepsilon^2$, because $z_g\ge rZ_b$.
It follows that
\[
 -t(\log M_t)_t\ge3/4.
\]
This is why the two physical sections cannot be identified: a departure
point arbitrarily close to its final target need not have a bounded
logarithmic normal derivative.

Since $B\Psi_b'/\Psi_b=1+O(B)$, conversion to $B$ yields
\[
 BM_B\le C\varepsilon^2(r^2+\xi_{\rm out}^2),\qquad
 -B(\log M_B)_B\ge2/3.
\]
For the complete map from $s$, $M=M_B(m_bA)m_bA_s$; the field is
fixed when $s$ is differentiated. Thus, exactly,
\begin{align*}
 sM&=(BM_B)K_A,\\
 -s(\log M)_s
 &=\{-B(\log M_B)_B\}K_A-(K_A-1)-s(\log K_A)_s\ge1/2.
\end{align*}
The unbounded multiplier has canceled. Choosing the fixed sections
small and then $r,\varepsilon$ small gives both estimates in
\eqref{eq:pc-outgoing-conditions} throughout $0<A<\eta_0/m_b$.
We now combine them with the fiber-exit estimates on the same returning
source interval.

For a fixed field, the returning source has the following partition.

\begin{lemma}[Source partition]
\label{lem:pc-source-partition}
On the actual source interval $I_{\rm ret}$ of
Lemma~\ref{lem:pc-collision-localization}, $A$ is a diffeomorphism from the
source coordinate $s$ onto an interval
$\mathcal A_{\rm actual}\subset(0,\eta_0)$. On the no-root side the
ordered open pieces are
\begin{equation}
\label{eq:pc-source-partition}
 \begin{aligned}
 \mathcal I_{\rm through}
   &=A^{-1}\!\left(\mathcal A_{\rm actual}
                    \cap(0,\eta_0/m_b)\right),\\
 \mathcal I_{\rm fiber}
   &=A^{-1}\!\left(\mathcal A_{\rm actual}
                    \cap(\eta_0/m_b,\eta_0)\right).
 \end{aligned}
\end{equation}
Either piece may be empty. There is at most one finite event value,
$A=\eta_0/m_b$, when it belongs to the actual source interval. On the
root side every source in $I_{\rm ret}$ exits through the fiber and
there is no internal event value. All counted cycles are in this
interval, and its artificial boundary $A=\eta_0$ contributes no counted
zero.
\end{lemma}
\begin{proof}
The construction \eqref{eq:pc-localized-source} and the positive section
conversion give $A_s>0$; \eqref{eq:pc-left-elasticity} gives its
quantitative form. On the no-root side the product equation
\eqref{eq:pc-through-multiplier} has $m_b>1$, so the only change of
exit is $m_bA=\eta_0$. Pulling back the two amplitude intervals gives
\eqref{eq:pc-source-partition}. At equality the trajectory meets the
right base and fiber at the same point, in finite time. If the threshold
is outside the actual amplitude image, it is omitted. The target-core
preimage remains an interval by scalar order and $F_y>0$, as in
Lemma~\ref{lem:sections}.

On the root side the scalar base tends to its first root in infinite
time, whereas $\dot\eta=b_n\eta$, $b_n\ge c>0$, takes every $A>0$ to
$\eta_0$ in finite time. Hence only the fiber exit occurs. The selected
boundary $A=0$ is excluded from the returning interval. Finally,
Lemma~\ref{lem:pc-collision-localization} places every counted cycle
strictly below $A=\eta_0/2$, so the upper artificial boundary introduces
no counted event zero. The ordering in Figure~\ref{fig:collision-passages} is
increasing $s$, or equivalently increasing $A$; it is opposite to the
original $u_{\rm up}$-order because $s_{u_{\rm up}}<0$.
\end{proof}

\paragraph{The complete return across the finite junction.}
On $I_{\rm ret}$, write $X(u)$ for the upper-section point and
$\phi_p$ for the fixed-field flow. The target hit
$\pi_z\phi_p(t,X(u))=\xi$ is transverse, with $F_p(\xi,y)\ge c_\xi>0$.
It defines one $C^2$ target map
$\widehat T_+(u)=E_b(\pi_y\phi_p(\tau(u),X(u)))$.
At $A=\eta_0/m_b$ the orbit meets the right base and fiber at their
common point in finite time; there $v>0$, $b_n\ge c>0$ and $A_s>0$.
Thus the two formulas are restrictions of this same map, and the
physical section diffeomorphisms give $M=\partial_s\widehat T_+>0$.
Its first two derivatives agree at the event, which is counted once.
A base root has infinite flight time and is never such a junction.
These are the flow and continuity hypotheses of
Lemma~\ref{lem:pc-event-gluing}; on the root side there is only fiber exit.

The through and fiber estimates hold on their whole respective returning
intervals. Fix the common positive curvature lower bound and the incoming
and short-map constants first, and choose $\delta_*>0$ satisfying both
inequalities in~\eqref{eq:pc-ratio-smallness} with these common bounds. The short
map constants can be fixed before the outgoing section: its limiting
energy map is independent of that section, and for each fixed section a
parameter shrinking bounds its $C^2$ error. Likewise the through size
constant is uniform in the fixed small energy rectangle; after fixing
the section, its final regular holonomy derivative is bounded by two.
Choose the fixed section ratio, then $\xi_{\rm out}$ sufficiently small,
and then $r$ so that the through coefficient
$C(r^2+\xi_{\rm out}^2)$ is below $\delta_*$. With this target fixed,
decrease $r$ further so that the fiber coefficient $C_\xi r^\alpha$
is below the same tolerance. Their curvature lower constants have a
common positive minimum, independently of $m_b$ and the root gap.
The localization lemma only decreases the final positive thresholds;
it changes none of these physical sections or coefficient intervals.

Taking the maximum size constant and minimum curvature constant, the
two passage estimates therefore give
\[
 sM\le\varepsilon^2\delta,\qquad
 -s(\log M)_s\ge c_M>0
\]
on both open pieces. The incoming and short-map estimates also hold on
$I_{\rm ret}$. These are precisely the hypotheses of
Lemma~\ref{lem:pc-event-gluing}; the continuity just proved includes the
possible junction. Its total actual multiplicity bound two therefore
holds on this interval, which contains every counted cycle.

\begin{proof}[Completion of Proposition~\ref{prop:positive-constant}]
For negative forcing, cover the coefficient set by the three intervals
$a\ge-2+\delta_a$, $a\le-2-\delta_a$, and $|a+2|<2\delta_a$.
For each fixed field choose one applicable interval. In either of the
first two, the preceding complete-return comparison gives the bound two.
In the third, Lemma~\ref{lem:pc-collision-localization} places every
counted cycle inside one actual source interval with $A<\eta_0/2$.
On the root side only the fiber exit remains. On the no-root side,
Lemma~\ref{lem:pc-source-partition} gives the through and fiber pieces
and at most one finite junction. Lemma~\ref{lem:pc-event-gluing}
therefore gives total multiplicity at most two on this common interval.

Every comparison counts the complete return from the single signed
source section to the common target. Fix the collision fiber before
shrinking its coefficient interval; this also fixes the gaps in the
separated-root regimes. Then fix the physical sections before the final
parameter bounds. A compact coefficient set requires only finitely many
such choices, whose minima give common positive thresholds
$r_0,\varepsilon_0$ and a common positive Hausdorff radius. Together
with the nonnegative-forcing comparison on the entire passing interval,
this proves the proposition.
\end{proof}
